%********************************************%
%*       Generated from PreTeXt source      *%
%*       on 2025-08-06T11:16:10-06:00       *%
%*   A recent stable commit (2022-07-01):   *%
%* 6c761d3dba23af92cba35001c852aac04ae99a5f *%
%*                                          *%
%*         https://pretextbook.org          *%
%*                                          *%
%********************************************%
\documentclass[oneside,10pt,]{book}
%% Custom Preamble Entries, early (use latex.preamble.early)
%% Default LaTeX packages
%%   1.  always employed (or nearly so) for some purpose, or
%%   2.  a stylewriter may assume their presence
\usepackage{geometry}
%% Some aspects of the preamble are conditional,
%% the LaTeX engine is one such determinant
\usepackage{ifthen}
%% etoolbox has a variety of modern conveniences
\usepackage{etoolbox}
\usepackage{ifxetex,ifluatex}
%% Raster graphics inclusion
\usepackage{graphicx}
%% Color support, xcolor package
%% Always loaded, for: add/delete text, author tools
%% Here, since tcolorbox loads tikz, and tikz loads xcolor
\PassOptionsToPackage{dvipsnames,svgnames,table}{xcolor}
\usepackage{xcolor}
%% begin: defined colors, via xcolor package, for styling
%% end: defined colors, via xcolor package, for styling
%% Colored boxes, and much more, though mostly styling
%% skins library provides "enhanced" skin, employing tikzpicture
%% boxes may be configured as "breakable" or "unbreakable"
%% "raster" controls grids of boxes, aka side-by-side
\usepackage{tcolorbox}
\tcbuselibrary{skins}
\tcbuselibrary{breakable}
\tcbuselibrary{raster}
%% We load some "stock" tcolorbox styles that we use a lot
%% Placement here is provisional, there will be some color work also
%% First, black on white, no border, transparent, but no assumption about titles
\tcbset{ bwminimalstyle/.style={size=minimal, boxrule=-0.3pt, frame empty,
colback=white, colbacktitle=white, coltitle=black, opacityfill=0.0} }
%% Second, bold title, run-in to text/paragraph/heading
%% Space afterwards will be controlled by environment,
%% independent of constructions of the tcb title
%% Places \blocktitlefont onto many block titles
\tcbset{ runintitlestyle/.style={fonttitle=\blocktitlefont\upshape\bfseries, attach title to upper} }
%% Spacing prior to each exercise, anywhere
\tcbset{ exercisespacingstyle/.style={before skip={1.5ex plus 0.5ex}} }
%% Spacing prior to each block
\tcbset{ blockspacingstyle/.style={before skip={2.0ex plus 0.5ex}} }
%% xparse allows the construction of more robust commands,
%% this is a necessity for isolating styling and behavior
%% The tcolorbox library of the same name loads the base library
\tcbuselibrary{xparse}
%% The tcolorbox library loads TikZ, its calc package is generally useful,
%% and is necessary for some smaller documents that use partial tcolor boxes
%% See:  https://github.com/PreTeXtBook/pretext/issues/1624
\usetikzlibrary{calc}
%% We use some more exotic tcolorbox keys to restore indentation to parboxes
\tcbuselibrary{hooks}
%% Save default paragraph indentation and parskip for use later, when adjusting parboxes
\newlength{\normalparindent}
\newlength{\normalparskip}
\AtBeginDocument{\setlength{\normalparindent}{\parindent}}
\AtBeginDocument{\setlength{\normalparskip}{\parskip}}
\newcommand{\setparstyle}{\setlength{\parindent}{\normalparindent}\setlength{\parskip}{\normalparskip}}%% Hyperref should be here, but likes to be loaded late
%%
%% Inline math delimiters, \(, \), need to be robust
%% 2016-01-31:  latexrelease.sty  supersedes  fixltx2e.sty
%% If  latexrelease.sty  exists, bugfix is in kernel
%% If not, bugfix is in  fixltx2e.sty
%% See:  https://tug.org/TUGboat/tb36-3/tb114ltnews22.pdf
%% and read "Fewer fragile commands" in distribution's  latexchanges.pdf
\IfFileExists{latexrelease.sty}{}{\usepackage{fixltx2e}}
%% Footnote counters and part/chapter counters are manipulated
%% April 2018:  chngcntr  commands now integrated into the kernel,
%% but circa 2018/2019 the package would still try to redefine them,
%% so we need to do the work of loading conditionally for old kernels.
%% From version 1.1a,  chngcntr  should detect defintions made by LaTeX kernel.
\ifdefined\counterwithin
\else
    \usepackage{chngcntr}
\fi
%% Text height identically 9 inches, text width varies on point size
%% See Bringhurst 2.1.1 on measure for recommendations
%% 75 characters per line (count spaces, punctuation) is target
%% which is the upper limit of Bringhurst's recommendations
\geometry{letterpaper,total={340pt,9.0in}}
%% Custom Page Layout Adjustments (use publisher page-geometry entry)
%% This LaTeX file may be compiled with pdflatex, xelatex, or lualatex executables
%% LuaTeX is not explicitly supported, but we do accept additions from knowledgeable users
%% The conditional below provides  pdflatex  specific configuration last
%% begin: engine-specific capabilities
\ifthenelse{\boolean{xetex} \or \boolean{luatex}}{%
%% begin: xelatex and lualatex-specific default configuration
\ifxetex\usepackage{xltxtra}\fi
%% realscripts is the only part of xltxtra relevant to lualatex 
\ifluatex\usepackage{realscripts}\fi
%% end:   xelatex and lualatex-specific default configuration
}{
%% begin: pdflatex-specific default configuration
%% We assume a PreTeXt XML source file may have Unicode characters
%% and so we ask LaTeX to parse a UTF-8 encoded file
%% This may work well for accented characters in Western language,
%% but not with Greek, Asian languages, etc.
%% When this is not good enough, switch to the  xelatex  engine
%% where Unicode is better supported (encouraged, even)
\usepackage[utf8]{inputenc}
%% end: pdflatex-specific default configuration
}
%% end:   engine-specific capabilities
%%
%% Fonts.  Conditional on LaTex engine employed.
%% Default Text Font: The Latin Modern fonts are
%% "enhanced versions of the [original TeX] Computer Modern fonts."
%% We use them as the default text font for PreTeXt output.
%% Default Monospace font: Inconsolata (aka zi4)
%% Sponsored by TUG: http://levien.com/type/myfonts/inconsolata.html
%% Loaded for documents with intentional objects requiring monospace
%% See package documentation for excellent instructions
%% fontspec will work universally if we use filename to locate OTF files
%% Loads the "upquote" package as needed, so we don't have to
%% Upright quotes might come from the  textcomp  package, which we also use
%% We employ the shapely \ell to match Google Font version
%% pdflatex: "varl" package option produces shapely \ell
%% pdflatex: "var0" package option produces plain zero (not used)
%% pdflatex: "varqu" package option produces best upright quotes
%% xelatex,lualatex: add OTF StylisticSet 1 for shapely \ell
%% xelatex,lualatex: add OTF StylisticSet 2 for plain zero (not used)
%% xelatex,lualatex: add OTF StylisticSet 3 for upright quotes
%%
%% Automatic Font Control
%% Portions of a document, are, or may, be affected by defined commands
%% These are perhaps more flexible when using  xelatex  rather than  pdflatex
%% The following definitions are meant to be re-defined in a style, using \renewcommand
%% They are scoped when employed (in a TeX group), and so should not be defined with an argument
\newcommand{\divisionfont}{\relax}
\newcommand{\blocktitlefont}{\relax}
\newcommand{\contentsfont}{\relax}
\newcommand{\pagefont}{\relax}
\newcommand{\tabularfont}{\relax}
\newcommand{\xreffont}{\relax}
\newcommand{\titlepagefont}{\relax}
%%
\ifthenelse{\boolean{xetex} \or \boolean{luatex}}{%
%% begin: font setup and configuration for use with xelatex
%% Generally, xelatex is necessary for non-Western fonts
%% fontspec package provides extensive control of system fonts,
%% meaning *.otf (OpenType), and apparently *.ttf (TrueType)
%% that live *outside* your TeX/MF tree, and are controlled by your *system*
%% (it is possible that a TeX distribution will place fonts in a system location)
%%
%% The fontspec package is the best vehicle for using different fonts in  xelatex
%% So we load it always, no matter what a publisher or style might want
%%
\usepackage{fontspec}
%%
%% begin: xelatex main font ("font-xelatex-main" template)
%% Latin Modern Roman is the default font for xelatex and so is loaded with a TU encoding
%% *in the format* so we can't touch it, only perhaps adjust it later
%% in one of two ways (then known by NFSS names such as "lmr")
%% (1) via NFSS with font family names such as "lmr" and "lmss"
%% (2) via fontspec with commands like \setmainfont{Latin Modern Roman}
%% The latter requires the font to be known at the system-level by its font name,
%% but will give access to OTF font features through optional arguments
%% https://tex.stackexchange.com/questions/470008/
%% where-and-how-does-fontspec-sty-specify-the-default-font-latin-modern-roman
%% http://tex.stackexchange.com/questions/115321
%% /how-to-optimize-latin-modern-font-with-xelatex
%%
%% end:   xelatex main font ("font-xelatex-main" template)
%% begin: xelatex mono font ("font-xelatex-mono" template)
%% (conditional on non-trivial uses being present in source)
\IfFontExistsTF{Inconsolatazi4-Regular.otf}{}{\GenericError{}{The font "Inconsolatazi4-Regular.otf" requested by PreTeXt output processed by the  xelatex  executable is not available.  Either a file cannot be located in default locations via a filename, or a font is not known by its name as part of your system.}{Consult the PreTeXt Guide for help with LaTeX fonts, or perhaps try using  pdflatex  as a test.}{}}
\IfFontExistsTF{Inconsolatazi4-Bold.otf}{}{\GenericError{}{The font "Inconsolatazi4-Bold.otf" requested by PreTeXt output processed by the  xelatex  executable is not available.  Either a file cannot be located in default locations via a filename, or a font is not known by its name as part of your system.}{Consult the PreTeXt Guide for help with LaTeX fonts, or perhaps try using  pdflatex  as a test.}{}}
\usepackage{zi4}
\setmonofont[BoldFont=Inconsolatazi4-Bold.otf,StylisticSet={1,3}]{Inconsolatazi4-Regular.otf}
%% end:   xelatex mono font ("font-xelatex-mono" template)
%% begin: xelatex font adjustments ("font-xelatex-style" template)
%% end:   xelatex font adjustments ("font-xelatex-style" template)
%%
%% Extensive support for other languages
\usepackage{polyglossia}
%% Set main/default language based on pretext/@xml:lang value
%% document language code is "en-US", US English
%% usmax variant has extra hypenation
\setmainlanguage[variant=usmax]{english}
%% Enable secondary languages based on discovery of @xml:lang values
%% Enable fonts/scripts based on discovery of @xml:lang values
%% Western languages should be ably covered by Latin Modern Roman
%% end:   font setup and configuration for use with xelatex
}{%
%% begin: font setup and configuration for use with pdflatex
%% begin: pdflatex main font ("font-pdflatex-main" template)
\usepackage{lmodern}
\usepackage[T1]{fontenc}
%% end:   pdflatex main font ("font-pdflatex-main" template)
%% begin: pdflatex mono font ("font-pdflatex-mono" template)
%% (conditional on non-trivial uses being present in source)
\usepackage[varqu,varl]{inconsolata}
%% end:   pdflatex mono font ("font-pdflatex-mono" template)
%% begin: pdflatex font adjustments ("font-pdflatex-style" template)
%% end:   pdflatex font adjustments ("font-pdflatex-style" template)
%% end:   font setup and configuration for use with pdflatex
}
%% Micromanage spacing, etc.  The named "microtype-options"
%% template may be employed to fine-tune package behavior
\usepackage{microtype}
%% Symbols, align environment, commutative diagrams, bracket-matrix
\usepackage{amsmath}
\usepackage{amscd}
\usepackage{amssymb}
%% allow page breaks within display mathematics anywhere
%% level 4 is maximally permissive
%% this is exactly the opposite of AMSmath package philosophy
%% there are per-display, and per-equation options to control this
%% split, aligned, gathered, and alignedat are not affected
\allowdisplaybreaks[4]
%% allow more columns to a matrix
%% can make this even bigger by overriding with  latex.preamble.late  processing option
\setcounter{MaxMatrixCols}{30}
%%
%%
%% Division Titles, and Page Headers/Footers
%% titlesec package, loading "titleps" package cooperatively
%% See code comments about the necessity and purpose of "explicit" option.
%% The "newparttoc" option causes a consistent entry for parts in the ToC 
%% file, but it is only effective if there is a \titleformat for \part.
%% "pagestyles" loads the  titleps  package cooperatively.
\usepackage[explicit, newparttoc, pagestyles]{titlesec}
%% The companion titletoc package for the ToC.
\usepackage{titletoc}
%% Fixes a bug with transition from chapters to appendices in a "book"
%% See generating XSL code for more details about necessity
\newtitlemark{\chaptertitlename}
%% begin: customizations of page styles via the modal "titleps-style" template
%% Designed to use commands from the LaTeX "titleps" package
%% Plain pages should have the same font for page numbers
\renewpagestyle{plain}{%
\setfoot{}{\pagefont\thepage}{}%
}%
%% Single pages as in default LaTeX
\renewpagestyle{headings}{%
\sethead{\pagefont\slshape\MakeUppercase{\ifthechapter{\chaptertitlename\space\thechapter.\space}{}\chaptertitle}}{}{\pagefont\thepage}%
}%
\pagestyle{headings}
%% end: customizations of page styles via the modal "titleps-style" template
%%
%% Create globally-available macros to be provided for style writers
%% These are redefined for each occurence of each division
\newcommand{\divisionnameptx}{\relax}%
\newcommand{\titleptx}{\relax}%
\newcommand{\subtitleptx}{\relax}%
\newcommand{\shortitleptx}{\relax}%
\newcommand{\authorsptx}{\relax}%
\newcommand{\epigraphptx}{\relax}%
%% Create environments for possible occurences of each division
%% Environment for a PTX "preface" at the level of a LaTeX "chapter"
\NewDocumentEnvironment{preface}{mmmmmmm}
{%
\renewcommand{\divisionnameptx}{#1}%
\renewcommand{\titleptx}{#2}%
\renewcommand{\subtitleptx}{#3}%
\renewcommand{\shortitleptx}{#4}%
\renewcommand{\authorsptx}{#5}%
\renewcommand{\epigraphptx}{#6}%
\chapter*{#2}%
\addcontentsline{toc}{chapter}{#4}
\label{#7}%
}{}%
%% Environment for a PTX "chapter" at the level of a LaTeX "chapter"
\NewDocumentEnvironment{chapterptx}{mmmmmmm}
{%
\renewcommand{\divisionnameptx}{#1}%
\renewcommand{\titleptx}{#2}%
\renewcommand{\subtitleptx}{#3}%
\renewcommand{\shortitleptx}{#4}%
\renewcommand{\authorsptx}{#5}%
\renewcommand{\epigraphptx}{#6}%
\chapter[{#4}]{#2}%
\label{#7}%
}{}%
%% Environment for a PTX "section" at the level of a LaTeX "section"
\NewDocumentEnvironment{sectionptx}{mmmmmmm}
{%
\renewcommand{\divisionnameptx}{#1}%
\renewcommand{\titleptx}{#2}%
\renewcommand{\subtitleptx}{#3}%
\renewcommand{\shortitleptx}{#4}%
\renewcommand{\authorsptx}{#5}%
\renewcommand{\epigraphptx}{#6}%
\section[{#4}]{#2}%
\label{#7}%
}{}%
%% Environment for a PTX "subsection" at the level of a LaTeX "subsection"
\NewDocumentEnvironment{subsectionptx}{mmmmmmm}
{%
\renewcommand{\divisionnameptx}{#1}%
\renewcommand{\titleptx}{#2}%
\renewcommand{\subtitleptx}{#3}%
\renewcommand{\shortitleptx}{#4}%
\renewcommand{\authorsptx}{#5}%
\renewcommand{\epigraphptx}{#6}%
\subsection[{#4}]{#2}%
\label{#7}%
}{}%
%% Environment for a PTX "exercises" at the level of a LaTeX "subsection"
\NewDocumentEnvironment{exercises-subsection}{mmmmmmm}
{%
\renewcommand{\divisionnameptx}{#1}%
\renewcommand{\titleptx}{#2}%
\renewcommand{\subtitleptx}{#3}%
\renewcommand{\shortitleptx}{#4}%
\renewcommand{\authorsptx}{#5}%
\renewcommand{\epigraphptx}{#6}%
\subsection[{#4}]{#2}%
\label{#7}%
}{}%
%% Environment for a PTX "exercises" at the level of a LaTeX "subsection"
\NewDocumentEnvironment{exercises-subsection-numberless}{mmmmmmm}
{%
\renewcommand{\divisionnameptx}{#1}%
\renewcommand{\titleptx}{#2}%
\renewcommand{\subtitleptx}{#3}%
\renewcommand{\shortitleptx}{#4}%
\renewcommand{\authorsptx}{#5}%
\renewcommand{\epigraphptx}{#6}%
\subsection*{#2}%
\addcontentsline{toc}{subsection}{#4}
\label{#7}%
}{}%
%% Environment for a PTX "solutions" at the level of a LaTeX "chapter"
\NewDocumentEnvironment{solutions-chapter}{mmmmmmm}
{%
\renewcommand{\divisionnameptx}{#1}%
\renewcommand{\titleptx}{#2}%
\renewcommand{\subtitleptx}{#3}%
\renewcommand{\shortitleptx}{#4}%
\renewcommand{\authorsptx}{#5}%
\renewcommand{\epigraphptx}{#6}%
\chapter[{#4}]{#2}%
\label{#7}%
}{}%
%% Environment for a PTX "solutions" at the level of a LaTeX "chapter"
\NewDocumentEnvironment{solutions-chapter-numberless}{mmmmmmm}
{%
\renewcommand{\divisionnameptx}{#1}%
\renewcommand{\titleptx}{#2}%
\renewcommand{\subtitleptx}{#3}%
\renewcommand{\shortitleptx}{#4}%
\renewcommand{\authorsptx}{#5}%
\renewcommand{\epigraphptx}{#6}%
\chapter*{#2}%
\addcontentsline{toc}{chapter}{#4}
\label{#7}%
}{}%
%%
%% Styles for six traditional LaTeX divisions
\titleformat{\part}[display]
{\divisionfont\Huge\bfseries\centering}{\divisionnameptx\space\thepart}{30pt}{\Huge#1}
[{\Large\centering\authorsptx}]
\titleformat{\chapter}[display]
{\divisionfont\huge\bfseries}{\divisionnameptx\space\thechapter}{20pt}{\Huge#1}
[{\Large\authorsptx}]
\titleformat{name=\chapter,numberless}[display]
{\divisionfont\huge\bfseries}{}{0pt}{#1}
[{\Large\authorsptx}]
\titlespacing*{\chapter}{0pt}{50pt}{40pt}
\titleformat{\section}[hang]
{\divisionfont\Large\bfseries}{\thesection}{1ex}{#1}
[{\large\authorsptx}]
\titleformat{name=\section,numberless}[block]
{\divisionfont\Large\bfseries}{}{0pt}{#1}
[{\large\authorsptx}]
\titlespacing*{\section}{0pt}{3.5ex plus 1ex minus .2ex}{2.3ex plus .2ex}
\titleformat{\subsection}[hang]
{\divisionfont\large\bfseries}{\thesubsection}{1ex}{#1}
[{\normalsize\authorsptx}]
\titleformat{name=\subsection,numberless}[block]
{\divisionfont\large\bfseries}{}{0pt}{#1}
[{\normalsize\authorsptx}]
\titlespacing*{\subsection}{0pt}{3.25ex plus 1ex minus .2ex}{1.5ex plus .2ex}
\titleformat{\subsubsection}[hang]
{\divisionfont\normalsize\bfseries}{\thesubsubsection}{1em}{#1}
[{\small\authorsptx}]
\titleformat{name=\subsubsection,numberless}[block]
{\divisionfont\normalsize\bfseries}{}{0pt}{#1}
[{\normalsize\authorsptx}]
\titlespacing*{\subsubsection}{0pt}{3.25ex plus 1ex minus .2ex}{1.5ex plus .2ex}
\titleformat{\paragraph}[hang]
{\divisionfont\normalsize\bfseries}{\theparagraph}{1em}{#1}
[{\small\authorsptx}]
\titleformat{name=\paragraph,numberless}[block]
{\divisionfont\normalsize\bfseries}{}{0pt}{#1}
[{\normalsize\authorsptx}]
\titlespacing*{\paragraph}{0pt}{3.25ex plus 1ex minus .2ex}{1.5em}
%%
%% Styles for five traditional LaTeX divisions
\titlecontents{part}%
[0pt]{\contentsmargin{0em}\addvspace{1pc}\contentsfont\bfseries}%
{\Large\thecontentslabel\enspace}{\Large}%
{}%
[\addvspace{.5pc}]%
\titlecontents{chapter}%
[0pt]{\contentsmargin{0em}\addvspace{1pc}\contentsfont\bfseries}%
{\large\thecontentslabel\enspace}{\large}%
{\hfill\bfseries\thecontentspage}%
[\addvspace{.5pc}]%
\dottedcontents{section}[3.8em]{\contentsfont}{2.3em}{1pc}%
\dottedcontents{subsection}[6.1em]{\contentsfont}{3.2em}{1pc}%
\dottedcontents{subsubsection}[9.3em]{\contentsfont}{4.3em}{1pc}%
%%
%% Begin: Semantic Macros
%% To preserve meaning in a LaTeX file
%%
%% \mono macro for content of "c", "cd", "tag", etc elements
%% Also used automatically in other constructions
%% Simply an alias for \texttt
%% Always defined, even if there is no need, or if a specific tt font is not loaded
\newcommand{\mono}[1]{\texttt{#1}}
%%
%% Following semantic macros are only defined here if their
%% use is required only in this specific document
%%
%% Used to markup initialisms, text or titles
\newcommand{\initialism}[1]{\textsc{\MakeLowercase{#1}}}
\DeclareRobustCommand{\initialismintitle}[1]{\texorpdfstring{#1}{#1}}
%% Used for warnings, typically bold and italic
\newcommand{\alert}[1]{\textbf{\textit{#1}}}
%% Used for inline definitions of terms
\newcommand{\terminology}[1]{\textbf{#1}}
%% Used for fillin answer blank in text
%% Relies on calc package, loaded via tcolorbox
%% Argument is intended number of characters of blank
%% Length may compress for output to fit in one line
\newlength{\fillinmaxwidth}
\newlength{\fillincontract}
\newlength{\charmaxwidth}\setlength{\charmaxwidth}{0.5em}
\newlength{\charminwidth}\setlength{\charminwidth}{0.1em}
\newlength{\fillinheight}
\newcommand{\fillintext}[1]{%
\setlength{\fillinmaxwidth}{#1\charmaxwidth}%
\setlength{\fillincontract}{#1\charminwidth}%
\setlength{\fillinheight}{\baselineskip}\addtolength{\fillinheight}{1.2pt}%
\strut\nobreak\leaders\vbox{\hrule width 0.3pt height 0.3pt \vskip -1.2pt}\hskip 1\fillinmaxwidth minus \fillincontract\nobreak\strut%
}
%% Style of a title on a list item, for ordered and unordered lists
%% Also "task" of exercise, PROJECT-LIKE, EXAMPLE-LIKE
\newcommand{\lititle}[1]{{\slshape#1}}
%% End: Semantic Macros
%% begin: environments for duplicates in solutions divisions
%% Solutions to inline exercises, style and environment
\tcbset{ inlinesolutionstyle/.style={bwminimalstyle, runintitlestyle, exercisespacingstyle, after title={\space}, breakable, before upper app={\setparstyle} } }
\newtcolorbox{inlinesolution}[4]{inlinesolutionstyle, title={\hyperref[#4]{#1~#2}\notblank{#3}{\space#3}{}}}
%% Solutions to division exercises, not in exercise group
%% Parameter #1 is type-name and is ignored
\tcbset{ divisionsolutionstyle/.style={bwminimalstyle, runintitlestyle, exercisespacingstyle, after title={\space}, breakable, before upper app={\setparstyle} } }
\newtcolorbox{divisionsolution}[4]{divisionsolutionstyle, title={\hyperlink{#4}{#2}.\notblank{#3}{\space#3}{}}}
%% Solutions to division exercises, in exercise group, no columns
%% Parameter #1 is type-name and is ignored
\tcbset{ divisionsolutionegstyle/.style={bwminimalstyle, runintitlestyle, exercisespacingstyle, after title={\space}, left skip=\egindent, breakable, before upper app={\setparstyle} } }
\newtcolorbox{divisionsolutioneg}[4]{divisionsolutionegstyle, title={\hyperlink{#4}{#2}.\notblank{#3}{\space#3}{}}}
%% Solutions to division exercises, in exercise group with columns
%% Parameter #1 is type-name and is ignored
\tcbset{ divisionsolutionegcolstyle/.style={bwminimalstyle, runintitlestyle,  exercisespacingstyle, after title={\space}, halign=flush left, unbreakable, before upper app={\setparstyle} } }
\newtcolorbox{divisionsolutionegcol}[4]{divisionsolutionegcolstyle, title={\hyperlink{#4}{#2}.\notblank{#3}{\space#3}{}}}
\tcbset{ projectsolutionstyle/.style={bwminimalstyle, runintitlestyle, exercisespacingstyle, after title={\space}, breakable, before upper app={\setparstyle} } }
\newtcolorbox{projectsolution}[4]{projectsolutionstyle, title={\hyperref[#4]{#1~#2}\notblank{#3}{\space#3}{}}}
%% Divisional exercises (and worksheet) as LaTeX environments
%% Third argument is option for extra workspace in worksheets
%% Hanging indent occupies a 5ex width slot prior to left margin
%% Experimentally this seems just barely sufficient for a bold "888."
%% Division exercises, not in exercise group
\tcbset{ divisionexercisestyle/.style={bwminimalstyle, runintitlestyle, exercisespacingstyle, left=5ex, breakable, before upper app={\setparstyle} } }
\newtcolorbox{divisionexercise}[4]{divisionexercisestyle, before title={\hspace*{-5ex}\makebox[5ex][l]{#1.}}, title={\notblank{#2}{#2}{}}, after title={\notblank{#2}{\space}{}}, phantom={\hypertarget{#4}{}}, after={\notblank{#3}{\par\rule{\workspacestrutwidth}{#3}\par\vfill}{\par}}}
%% Division exercises, in exercise group, no columns
\tcbset{ divisionexerciseegstyle/.style={bwminimalstyle, runintitlestyle, exercisespacingstyle, left=5ex, left skip=\egindent, breakable, before upper app={\setparstyle} } }
\newtcolorbox{divisionexerciseeg}[4]{divisionexerciseegstyle, before title={\hspace*{-5ex}\makebox[5ex][l]{#1.}}, title={\notblank{#2}{#2}{}}, after title={\notblank{#2}{\space}{}}, phantom={\hypertarget{#4}{}}, after={\notblank{#3}{\par\rule{\workspacestrutwidth}{#3}\par\vfill}{\par}}}
%% Division exercises, in exercise group with columns
\tcbset{ divisionexerciseegcolstyle/.style={bwminimalstyle, runintitlestyle, exercisespacingstyle, left=5ex, halign=flush left, unbreakable, before upper app={\setparstyle} } }
\newtcolorbox{divisionexerciseegcol}[4]{divisionexerciseegcolstyle, before title={\hspace*{-5ex}\makebox[5ex][l]{#1.}}, title={\notblank{#2}{#2}{}}, after title={\notblank{#2}{\space}{}}, phantom={\hypertarget{#4}{}}, after upper={\notblank{#3}{\par\rule{\workspacestrutwidth}{#3}\par\vfill}{\par}}}
%% Equation Numbering
%% Controlled by  numbering.equations.level  processing parameter
%% No adjustment here implies document-wide numbering
\numberwithin{equation}{section}
%% "tcolorbox" environment for a single image, occupying entire \linewidth
%% arguments are left-margin, width, right-margin, as multiples of
%% \linewidth, and are guaranteed to be positive and sum to 1.0
\tcbset{ imagestyle/.style={bwminimalstyle} }
\NewTColorBox{tcbimage}{mmm}{imagestyle,left skip=#1\linewidth,width=#2\linewidth}
%% Wrapper environment for tcbimage environment with a fourth argument
%% Fourth argument, if nonempty, is a vertical space adjustment
%% and implies image will be preceded by \leavevmode\nopagebreak
%% Intended use is for alignment with a list marker
\NewDocumentEnvironment{image}{mmmm}{\notblank{#4}{\leavevmode\nopagebreak\vspace{#4}}{}\begin{tcbimage}{#1}{#2}{#3}}{\end{tcbimage}%
}%% For improved tables
\usepackage{array}
%% Some extra height on each row is desirable, especially with horizontal rules
%% Increment determined experimentally
\setlength{\extrarowheight}{0.2ex}
%% Define variable thickness horizontal rules, full and partial
%% Thicknesses are 0.03, 0.05, 0.08 in the  booktabs  package
\newcommand{\hrulethin}  {\noalign{\hrule height 0.04em}}
\newcommand{\hrulemedium}{\noalign{\hrule height 0.07em}}
\newcommand{\hrulethick} {\noalign{\hrule height 0.11em}}
%% We preserve a copy of the \setlength package before other
%% packages (extpfeil) get a chance to load packages that redefine it
\let\oldsetlength\setlength
\newlength{\Oldarrayrulewidth}
\newcommand{\crulethin}[1]%
{\noalign{\global\oldsetlength{\Oldarrayrulewidth}{\arrayrulewidth}}%
\noalign{\global\oldsetlength{\arrayrulewidth}{0.04em}}\cline{#1}%
\noalign{\global\oldsetlength{\arrayrulewidth}{\Oldarrayrulewidth}}}%
\newcommand{\crulemedium}[1]%
{\noalign{\global\oldsetlength{\Oldarrayrulewidth}{\arrayrulewidth}}%
\noalign{\global\oldsetlength{\arrayrulewidth}{0.07em}}\cline{#1}%
\noalign{\global\oldsetlength{\arrayrulewidth}{\Oldarrayrulewidth}}}
\newcommand{\crulethick}[1]%
{\noalign{\global\oldsetlength{\Oldarrayrulewidth}{\arrayrulewidth}}%
\noalign{\global\oldsetlength{\arrayrulewidth}{0.11em}}\cline{#1}%
\noalign{\global\oldsetlength{\arrayrulewidth}{\Oldarrayrulewidth}}}
%% Single letter column specifiers defined via array package
\newcolumntype{A}{!{\vrule width 0.04em}}
\newcolumntype{B}{!{\vrule width 0.07em}}
\newcolumntype{C}{!{\vrule width 0.11em}}
%% tcolorbox to place tabular outside of a sidebyside
\tcbset{ tabularboxstyle/.style={bwminimalstyle,} }
\newtcolorbox{tabularbox}[3]{tabularboxstyle, left skip=#1\linewidth, width=#2\linewidth,}
%% Footnote Numbering
%% Specified by numbering.footnotes.level
%% Undo counter reset by chapter for a book
\counterwithout{footnote}{chapter}
\counterwithin*{footnote}{section}
%% Fancy Verbatim for consoles, preformatted, code display, literate programming
\usepackage{fancyvrb}
%% Pre-formatted text, a peer of paragraphs
\DefineVerbatimEnvironment{preformatted}{Verbatim}{}
%% More flexible list management, esp. for references
%% But also for specifying labels (i.e. custom order) on nested lists
\usepackage[inline]{enumitem}
%% Description lists as tcolorbox sidebyside
%% "dli" short for "description list item"
\newlength{\dlititlewidth}
\newlength{\dlimaxnarrowtitle}\setlength{\dlimaxnarrowtitle}{11ex}
\newlength{\dlimaxmediumtitle}\setlength{\dlimaxmediumtitle}{18ex}
\tcbset{ dlistyle/.style={sidebyside, sidebyside align=top seam, lower separated=false, bwminimalstyle, bottomtitle=0.75ex, after skip=1.5ex, boxsep=0pt, left=0pt, right=0pt, top=0pt, bottom=0pt, before lower app={\setparstyle\noindent}} }
\tcbset{ dlinarrowstyle/.style={dlistyle, lefthand width=\dlimaxnarrowtitle, sidebyside gap=1ex, halign=flush left, righttitle=10ex} }
\tcbset{ dlimediumstyle/.style={dlistyle, lefthand width=\dlimaxmediumtitle, sidebyside gap=4ex, halign=flush right} }
\NewDocumentEnvironment{descriptionlist}{}{\par\vspace*{1.5ex}}{\par\vspace*{1.5ex}}%
%% begin enviroment has an if/then to open the tcolorbox
\NewDocumentEnvironment{dlinarrow}{mm}{%
\settowidth{\dlititlewidth}{{\textbf{#1}}}%
\ifthenelse{\dlititlewidth > \dlimaxnarrowtitle}%
{\begin{tcolorbox}[title={\textbf{#1}}, phantom={\hypertarget{#2}{}}, dlinarrowstyle]\tcblower}%
{\begin{tcolorbox}[dlinarrowstyle, phantom={\hypertarget{#2}{}}]\textbf{#1}\tcblower}%
}%
{\end{tcolorbox}}%
%% medium option is simpler
\NewDocumentEnvironment{dlimedium}{mm}%
{\begin{tcolorbox}[dlimediumstyle, phantom={\hypertarget{#2}{}}]\textbf{#1}\tcblower}%
{\end{tcolorbox}}%
%% Indented groups of "exercise" within an "exercises" division
%% Lengths control the indentation (always) and gaps (multi-column)
\newlength{\egindent}\setlength{\egindent}{0.05\linewidth}
\newlength{\exggap}\setlength{\exggap}{0.05\linewidth}
%% Thin "xparse" environments will represent the entire exercise
%% group, in the case when it does not hold multiple columns.
\NewDocumentEnvironment{exercisegroup}{}
{}{}
%% An exercise group with multiple columns is a tcbraster.
%% If the contained exercises are explicitly unbreakable,
%% the raster should break at rows for page breaks.
%% The number of columns is a parameter, passed to tcbraster.
\tcbset{ exgroupcolstyle/.style={raster equal height=rows, raster left skip=\egindent, raster column skip=\exggap} }
\NewDocumentEnvironment{exercisegroupcol}{m}
{\begin{tcbraster}[exgroupcolstyle,raster columns=#1]}{\end{tcbraster}}
%% hyperref driver does not need to be specified, it will be detected
%% Footnote marks in tcolorbox have broken linking under
%% hyperref, so it is necessary to turn off all linking
%% It *must* be given as a package option, not with \hypersetup
\usepackage[hyperfootnotes=false]{hyperref}
%% configure hyperref's  \href{}{}  and  \nolinkurl  to match listings' inline verbatim
\renewcommand\UrlFont{\small\ttfamily}
%% Hyperlinking active in electronic PDFs, all links without surrounding boxes and blue
\hypersetup{colorlinks=true,linkcolor=blue,citecolor=blue,filecolor=blue,urlcolor=blue}
%% Less-clever names for hyperlinks are more reliable, *especially* for structural parts
%% See comments in the code to learn more about the importance of this setting
\hypersetup{hypertexnames=false}
%%The  hypertexnames  setting then confuses the hyperlinking from the index
%%This patch resolves the incorrect links, see code for StackExchange post.
\makeatletter
\patchcmd\Hy@EveryPageBoxHook{\Hy@EveryPageAnchor}{\Hy@hypertexnamestrue\Hy@EveryPageAnchor}{}{\fail}
\makeatother
\hypersetup{pdftitle={Integrating WeBWorK into Textbooks}}
%% If you manually remove hyperref, leave in this next command
%% This will allow LaTeX compilation, employing this no-op command
\providecommand\phantomsection{}
%% Division Numbering: Chapters, Sections, Subsections, etc
%% Division numbers may be turned off at some level ("depth")
%% A section *always* has depth 1, contrary to us counting from the document root
%% The latex default is 3.  If a larger number is present here, then
%% removing this command may make some cross-references ambiguous
%% The precursor variable $numbering-maxlevel is checked for consistency in the common XSL file
\setcounter{secnumdepth}{3}
%%
%%
%% A faux tcolorbox whose only purpose is to provide common numbering
%% facilities for most blocks (possibly not projects, 2D displays)
%% Controlled by  numbering.theorems.level  processing parameter
\newtcolorbox[auto counter, number within=section]{block}{}
%%
%% This document is set to number PROJECT-LIKE on a separate numbering scheme
%% So, a faux tcolorbox whose only purpose is to provide this numbering
%% Controlled by  numbering.projects.level  processing parameter
\newtcolorbox[auto counter, number within=section]{project-distinct}{}
%% A faux tcolorbox whose only purpose is to provide common numbering
%% facilities for 2D displays which are subnumbered as part of a "sidebyside"
\makeatletter
\newtcolorbox[auto counter, number within=tcb@cnt@block, number freestyle={\noexpand\thetcb@cnt@block(\noexpand\alph{\tcbcounter})}]{subdisplay}{}
\makeatother
%%
%% tcolorbox, with styles, for THEOREM-LIKE
%%
%% theorem: fairly simple numbered block/structure
\tcbset{ theoremstyle/.style={bwminimalstyle, runintitlestyle, blockspacingstyle, after title={\space}, before upper app={\setparstyle}, } }
\newtcolorbox[use counter from=block]{theorem}[4]{title={{#1~\thetcbcounter\notblank{#2#3}{\space}{}\notblank{#2}{\space#2}{}\notblank{#3}{\space(#3)}{}}}, phantomlabel={#4}, breakable, after={\par}, fontupper=\itshape, theoremstyle, }
%% corollary: fairly simple numbered block/structure
\tcbset{ corollarystyle/.style={bwminimalstyle, runintitlestyle, blockspacingstyle, after title={\space}, before upper app={\setparstyle}, } }
\newtcolorbox[use counter from=block]{corollary}[4]{title={{#1~\thetcbcounter\notblank{#2#3}{\space}{}\notblank{#2}{\space#2}{}\notblank{#3}{\space(#3)}{}}}, phantomlabel={#4}, breakable, after={\par}, fontupper=\itshape, corollarystyle, }
%%
%% tcolorbox, with styles, for PROOF-LIKE
%%
%% proof: title is a replacement
\tcbset{ proofstyle/.style={bwminimalstyle, fonttitle=\blocktitlefont\itshape, attach title to upper, after title={\space}, after upper={\space\space\hspace*{\stretch{1}}\(\blacksquare\)},
} }
\newtcolorbox{proof}[3]{title={\notblank{#2}{#2}{#1.}}, phantom={\hypertarget{#3}{}}, breakable, after={\par}, proofstyle, before upper app={\setparstyle} }
%%
%% tcolorbox, with styles, for DEFINITION-LIKE
%%
%% definition: fairly simple numbered block/structure
\tcbset{ definitionstyle/.style={bwminimalstyle, runintitlestyle, blockspacingstyle, after title={\space}, after upper={\space\space\hspace*{\stretch{1}}\(\lozenge\)}, before upper app={\setparstyle}, } }
\newtcolorbox[use counter from=block]{definition}[3]{title={{#1~\thetcbcounter\notblank{#2}{\space\space#2}{}}}, phantomlabel={#3}, breakable, after={\par}, definitionstyle, }
%%
%% tcolorbox, with styles, for REMARK-LIKE
%%
%% warning: fairly simple numbered block/structure
\tcbset{ warningstyle/.style={bwminimalstyle, runintitlestyle, blockspacingstyle, after title={\space}, before upper app={\setparstyle}, } }
\newtcolorbox[use counter from=block]{warning}[3]{title={{#1~\thetcbcounter\notblank{#2}{\space\space#2}{}}}, phantomlabel={#3}, breakable, after={\par}, warningstyle, }
%%
%% tcolorbox, with styles, for inline exercises
%%
%% inlineexercise: fairly simple numbered block/structure
\tcbset{ inlineexercisestyle/.style={bwminimalstyle, runintitlestyle, blockspacingstyle, after title={\space}, before upper app={\setparstyle}, } }
\newtcolorbox[use counter from=block]{inlineexercise}[3]{title={{#1~\thetcbcounter\notblank{#2}{\space\space#2}{}}}, phantomlabel={#3}, breakable, after={\par}, inlineexercisestyle, }
%%
%% tcolorbox, with styles, for PROJECT-LIKE
%%
%% project: fairly simple numbered block/structure
\tcbset{ projectstyle/.style={bwminimalstyle, runintitlestyle, blockspacingstyle, after title={\space}, before upper app={\setparstyle}, } }
\newtcolorbox[use counter from=project-distinct]{project}[3]{title={{#1~\thetcbcounter\notblank{#2}{\space\space#2}{}}}, phantomlabel={#3}, breakable, after={\par}, projectstyle, }
%%
%% xparse environments for introductions and conclusions of divisions
%%
%% introduction: in a structured division
\NewDocumentEnvironment{introduction}{m}
{\notblank{#1}{\noindent\textbf{#1}\space}{}}{\par\medskip}
%% Graphics Preamble Entries
\usepackage{pgfplots}
\pgfplotsset{
    every axis/.append style={
        axis lines=middle,
        xlabel={\(x\)},
        ylabel={\(y\)},
        grid = both,
    }
}
%% If tikz has been loaded, replace ampersand with \amp macro
\ifdefined\tikzset
    \tikzset{ampersand replacement = \amp}
\fi
%% Custom Preamble Entries, late (use latex.preamble.late)
%% extpfeil package for certain extensible arrows,
%% as also provided by MathJax extension of the same name
%% NB: this package loads mtools, which loads calc, which redefines
%%     \setlength, so it can be removed if it seems to be in the 
%%     way and your math does not use:
%%     
%%     \xtwoheadrightarrow, \xtwoheadleftarrow, \xmapsto, \xlongequal, \xtofrom
%%     
%%     we have had to be extra careful with variable thickness
%%     lines in tables, and so also load this package late
\usepackage{extpfeil}
%% Begin: Author-provided macros
%% (From  docinfo/macros  element)
%% Plus three from PTX for XML characters
\newcommand{\definiteintegral}[4]{\int_{#1}^{#2}\,#3\,d#4}
\newcommand{\indefiniteintegral}[2]{\int#1\,d#2} 
\newcommand{\lt}{<}
\newcommand{\gt}{>}
\newcommand{\amp}{&}
%% End: Author-provided macros
\begin{document}
%% bottom alignment is explicit, since it normally depends on oneside, twoside
\raggedbottom
\frontmatter
%% begin: half-title
\thispagestyle{empty}
{\titlepagefont\centering
\vspace*{0.28\textheight}
{\Huge Integrating WeBWorK into Textbooks}\\[2\baselineskip]
{\LARGE Sample Exercises}\\
}
\clearpage
%% end:   half-title
%% begin: title page
%% Inspired by Peter Wilson's "titleDB" in "titlepages" CTAN package
\thispagestyle{empty}
{\titlepagefont\centering
\vspace*{0.14\textheight}
%% Target for xref to top-level element is ToC
\addtocontents{toc}{\protect\hypertarget{sample-ww-chapter}{}}
{\Huge Integrating WeBWorK into Textbooks}\\[\baselineskip]
{\LARGE Sample Exercises}\\[3\baselineskip]
{\Large Robert Beezer}\\[0.5\baselineskip]
{\Large University of Puget Sound\\
Tacoma, Washington, USA}\\[3\baselineskip]
{\Large Alex Jordan}\\[0.5\baselineskip]
{\Large Portland Community College\\
Portland, Oregon, USA}\\[3\baselineskip]
{\Large August 6, 2025}\\}
\clearpage
%% end:   title page
%% begin: copyright-page
\thispagestyle{empty}
\vspace*{\stretch{2}}
\vspace*{\stretch{1}}
\null\clearpage
%% end:   copyright-page
%
%
\typeout{************************************************}
\typeout{Preface  Preface}
\typeout{************************************************}
%
\begin{preface}{Preface}{Preface}{}{Preface}{}{}{frontmatter-3}
This is a sample book to demonstrate integrating WeBWorK homework problems into content authored with PreTeXt.  It is structured as a book whose chapters are again structured with sections (no introductions or conclusions) so that it fits Ruestone's chaper\slash{}subchapter model and can be used for testing on Runestone Academy.%
\par
This work was funded with grants from \href{http://openoregon.org/}{OpenOregon}\footnote{\nolinkurl{openoregon.org/}\label{frontmatter-3-2-2}} and the \href{http://pugetsound.edu}{University of Puget Sound}\footnote{\nolinkurl{pugetsound.edu}\label{frontmatter-3-2-4}}.%
\par
Rob Beezer (University of Puget Sound) and Alex Jordan (Portland Community College) worked on the PreTeXt enhancements that make this possible. Mike Gage (Univeristy of Rochester), Geoff Goehle (Western Carolina University), and Alex Jordan made this possible by enhancing the WeBWorK end, and generally maintaining WeBWorK software.%
\par
This book assumes a mild familiarity with both \href{https://pretextbook.org/}{PreTeXt}\footnote{\nolinkurl{pretextbook.org/}\label{frontmatter-3-4-2}} and \href{http://webwork.maa.org/wiki/Main_Page}{WeBWorK}\footnote{\nolinkurl{webwork.maa.org/wiki/Main_Page}\label{frontmatter-3-4-4}}. For more information about either project, follow the links.%
\end{preface}
%% begin: table of contents
%% Adjust Table of Contents
\setcounter{tocdepth}{1}
\renewcommand*\contentsname{Contents}
\tableofcontents
%% end:   table of contents
\mainmatter
%
%
\typeout{************************************************}
\typeout{Chapter 1 Instructive Examples}
\typeout{************************************************}
%
\begin{chapterptx}{Chapter}{Instructive Examples}{}{Instructive Examples}{}{}{sample-ww-chapter-4}
\renewcommand*{\chaptername}{Chapter}
%
%
\typeout{************************************************}
\typeout{Section 1.1 Arithmetic}
\typeout{************************************************}
%
\begin{sectionptx}{Section}{Arithmetic}{}{Arithmetic}{}{}{sample-ww-chapter-4-2}
Some questions with quantitative answers.%
\begin{inlineexercise}{Checkpoint}{Adding Single-Digit Integers.}{integer-addition}%
A simple, but functional example to begin with. If you are just learning how to add, you can test yourself here.%
\par
Compute the sum of \(6\) and \(1\text{:}\)%
\par
\(6 + 1 =\) \fillintext{5}%
\par
That was a simple problem.  Let's move on.%
\par\smallskip%
\noindent\textbf{\blocktitlefont Answer}.\hypertarget{integer-addition-3}{}\quad{}\(7\)%
\par\smallskip%
\noindent\textbf{\blocktitlefont Solution}.\hypertarget{integer-addition-4}{}\quad{}\(6 + 1 = {7}\text{.}\)%
\end{inlineexercise}%
If you are familiar with WeBWorK, then it may be a surprise to you to be interacting with a WeBWorK problem this way, without having logged in to WeBWorK.%
\begin{inlineexercise}{Checkpoint}{Declaring a Problem Seed.}{integer-addition-with-seed}%
You can also declare a \mono{seed} to specify a version of any problem that has randomization. Here is the same problem (``copied'' in the PreTeXt source), but with a \mono{seed} specified.%
\par
Compute the sum of \(5\) and \(8\text{:}\)%
\par
\(5 + 8 =\) \fillintext{5}%
\par\smallskip%
\noindent\textbf{\blocktitlefont Answer}.\hypertarget{integer-addition-with-seed-3}{}\quad{}\(13\)%
\par\smallskip%
\noindent\textbf{\blocktitlefont Solution}.\hypertarget{integer-addition-with-seed-4}{}\quad{}\(5 + 8 = {13}\text{.}\)%
\end{inlineexercise}%
\begin{inlineexercise}{Checkpoint}{Controlling Randomness.}{integer-addition-with-control-seed}%
You can code your problem with randomization, but then use a specific \mono{seed} and WeBWorK's \mono{\$envir\{problemSeed\}} to override that randomization for the purposes of the version that will appear in HTML and print output.%
\par
Compute the sum of \(1\) and \(2\text{:}\)%
\par
\(1 + 2 =\) \fillintext{5}%
\par\smallskip%
\noindent\textbf{\blocktitlefont Answer}.\hypertarget{integer-addition-with-control-seed-3}{}\quad{}\(3\)%
\par\smallskip%
\noindent\textbf{\blocktitlefont Solution}.\hypertarget{integer-addition-with-control-seed-4}{}\quad{}\(1 + 2 = {3}\text{.}\)%
\end{inlineexercise}%
\begin{inlineexercise}{Checkpoint}{Special Answer Checking.}{special-answer-checking}%
One of the strengths of WeBWorK is its ability to give intelligent feedback for incorrect answers.%
\begin{itemize}[label=\textbullet]
\item{}There is general feedback for when the student's answer is in an entirely different ballpark from the correct answer. Try entering something like ``\mono{y}''.%
\item{}There is general feedback for when the student's answer is not in the right form. Try entering ``\mono{x\textasciicircum{}2*x\textasciicircum{}3}'', which, right or wrong, is unsimplified.%
\item{}And problems can be written to detect and respond to common mistakes. Try entering an answer where you multiply the two exponents (instead of adding them, which would be correct.)%
\end{itemize}
%
\par
Use the properties of exponents to simplify \({x^{5}x^{3}}\text{.}\)%
\par
\({x^{5}x^{3}}=\) \fillintext{8}%
\par\smallskip%
\noindent\textbf{\blocktitlefont Answer}.\hypertarget{special-answer-checking-3}{}\quad{}\(x^{8}\)%
\par\smallskip%
\noindent\textbf{\blocktitlefont Solution}.\hypertarget{special-answer-checking-4}{}\quad{}We \emph{add} the exponents as follows, while including a gratuitous reference to the quadratic formula:%
\par
%
\begin{equation*}
\begin{aligned}
{x^{5}x^{3}}\amp =x^{5+3}\amp \text{Theorem 1.2.1}\\
\amp =x^{8}
\end{aligned}
\end{equation*}
%
\end{inlineexercise}%
\begin{inlineexercise}{Checkpoint}{Using Hints.}{using-hints}%
Hints can be inserted into exercises. Whether a hint is visible in the HTML depends on the value of \mono{\$showHint} in \mono{PGcourse.pl} in the WeBWorK course that is hosting these exercises. How the hint is displayed in the pdf output is controlled in the usual way that an PTX hint is controlled.%
\par
Simplify the expression \({\sqrt{48}}\text{.}\)%
\par
\({\sqrt{48}}=\) \fillintext{10}%
\par\smallskip%
\noindent\textbf{\blocktitlefont Hint}.\hypertarget{using-hints-3}{}\quad{}Factor the number inside the radical.%
\par\smallskip%
\noindent\textbf{\blocktitlefont Answer}.\hypertarget{using-hints-4}{}\quad{}\(4\sqrt{3}\)%
\par\smallskip%
\noindent\textbf{\blocktitlefont Solution}.\hypertarget{using-hints-5}{}\quad{}%
\begin{equation*}
{\sqrt{48}}=\sqrt{4^2\cdot 3}=4\sqrt{3}
\end{equation*}
%
\end{inlineexercise}%
\begin{inlineexercise}{Checkpoint}{No Randomization.}{no-random}%
This problem has no randomization at all, not even if it were exported to be part of a problem set on a WeBWorK server. As such, it really doesn't need any lines of Perl code in its setup, so you have the option of skipping that part of the authoring process.%
\par
\(1+2=\) \fillintext{5}%
\par
For more about problems that do not require any randomization, see the \href{https://pretextbook.org/documentation.html}{PTX Author Guide}\footnotemark{}.%
\par\smallskip%
\noindent\textbf{\blocktitlefont Answer}.\hypertarget{no-random-3}{}\quad{}\(3\)%
\end{inlineexercise}%
\footnotetext[1]{\nolinkurl{pretextbook.org/documentation.html}\label{no-random-2-3-2}}%
\begin{project}{Project}{Inside a \mono{<project>}.}{copy-webwork-add-numbers}%
If you like, you can have a WeBWorK inside a PROJECT-LIKE block. Just like with an \mono{<exercise>}, it can be preceded with an optional \mono{<introduction>} and followed by an optional \mono{<conclusion>}.%
\par
Compute the sum of \(6\) and \(5\text{:}\)%
\par
\(6 + 5 =\) \fillintext{5}%
\par\smallskip%
\noindent\textbf{\blocktitlefont Answer}.\hypertarget{copy-webwork-add-numbers-3}{}\quad{}\(11\)%
\par\smallskip%
\noindent\textbf{\blocktitlefont Solution}.\hypertarget{copy-webwork-add-numbers-4}{}\quad{}\(6 + 5 = {11}\text{.}\)%
\end{project}%
\end{sectionptx}
%
%
\typeout{************************************************}
\typeout{Section 1.2 The Quadratic Formula}
\typeout{************************************************}
%
\begin{sectionptx}{Section}{The Quadratic Formula}{}{The Quadratic Formula}{}{}{section-the-quadratic-formula}
In the previous section, we saw relatively simple WeBWorK questions. This section demonstrates how even very complicated WeBWorK problems can still behave well.%
\par
Here is a theorem that gives us a formula for the solutions of a second-degree polynomial equation. Note later how the WeBWorK problem references the theorem by its number. This seemingly minor detail demonstrates the degree to which WeBWorK and PreTeXt have been integrated.%
\begin{theorem}{Theorem}{Quadratic Formula.}{}{theorem-quadratic-formula}%
Given the second-degree polynomial equation \(ax^2 + bx + c = 0\), where \(a\neq0\), solutions are given by%
\begin{equation*}
x = \frac{-b\pm\sqrt{b^2-4ac}}{2a}.
\end{equation*}
%
\end{theorem}
\begin{proof}{Proof}{}{theorem-quadratic-formula-3}
%
\begin{align*}
ax^2 + bx + c &= 0\\
ax^2 + bx &= -c\\
4ax^2 + 4bx &= -4c\\
4ax^2 + 4bx + b^2 &= b^2 - 4ac\\
(2ax + b)^2 &= b^2 - 4ac\\
2ax + b &=\pm\sqrt{b^2 - 4ac}\\
2ax &=-b\pm\sqrt{b^2 - 4ac}\\
x &=\frac{-b\pm\sqrt{b^2 - 4ac}}{2a}
\end{align*}
%
\end{proof}
\begin{inlineexercise}{Checkpoint}{Solving Quadratic Equations.}{quadratic-equation}%
Consider the quadratic equation \({5x^{2}-6x-8} = 0\text{.}\)%
\begin{enumerate}[font=\bfseries,label=(\alph*),ref=\alph*]%
\item{}\lititle{Identify Coefficients.}\par%
Identify the coefficients for the quadratic equation using the standard form from Theorem~1.2.1.%
\par
\(a=\) \fillintext{4}, \(b=\) \fillintext{4}, \(c=\) \fillintext{4}%
\par\smallskip%
\noindent\textbf{\blocktitlefont Answer 1}.\hypertarget{quadratic-equation-3-3}{}\quad{}\(5\)%
\par\smallskip%
\noindent\textbf{\blocktitlefont Answer 2}.\hypertarget{quadratic-equation-3-4}{}\quad{}\(-6\)%
\par\smallskip%
\noindent\textbf{\blocktitlefont Answer 3}.\hypertarget{quadratic-equation-3-5}{}\quad{}\(-8\)%
\par\smallskip%
\noindent\textbf{\blocktitlefont Solution}.\hypertarget{quadratic-equation-3-6}{}\quad{}Take the coefficient of \(x^2\) for the value of \(a\text{,}\) the coefficient of \(x\) for \(b\text{,}\) and the constant for \(c\text{.}\) \emph{In this case}, they are \(a = {5}\text{,}\) \(b = {-6}\text{,}\) \(c = {-8}\text{.}\)%
\item{}\lititle{Use the Quadratic Formula.}\par%
Using the quadratic formula, solve the equation.%
\par\smallskip%
\noindent\textbf{\blocktitlefont Answer}.\hypertarget{quadratic-equation-4-3}{}\quad{}\(\left\{2,\frac{-4}{5}\right\}\)%
\par\smallskip%
\noindent\textbf{\blocktitlefont Solution}.\hypertarget{quadratic-equation-4-4}{}\quad{}Recall that the quadratic formula is given in Theorem~1.2.1.%
\par
You already identified \(a = {5}\text{,}\) \(b = {-6}\text{,}\) and \(c = {-8}\text{,}\) so the results are:%
\par
%
\begin{equation*}
x = {\frac{-\left(-6\right)+\sqrt{\left(-6\right)^{2}-4\cdot 5\cdot \left(-8\right)}}{2\cdot 5}} = {2}
\end{equation*}
%
\par
or%
\par
%
\begin{equation*}
x = {\frac{-\left(-6\right)-\sqrt{\left(-6\right)^{2}-4\cdot 5\cdot \left(-8\right)}}{2\cdot 5}} = {-{\frac{4}{5}}}
\end{equation*}
%
\end{enumerate}%
This conclusion is just here for testing.%
\end{inlineexercise}%
\begin{inlineexercise}{Checkpoint}{Nested tasks.}{nested-tasks}%
This exercise tests that nested tasks work.%
 \par
Consider the quadratic equation \({6x^{2}-31x-30} = 0\text{.}\)%
\begin{enumerate}[font=\bfseries,label=(\alph*),ref=\alph*]%
\item{}\lititle{Identify Coefficients.}\par%
Identify the coefficients for the quadratic equation using the standard form from Theorem~1.2.1.%
\begin{enumerate}[font=\bfseries,label=(\roman*),ref=\theenumi.\roman*]%
\item{}\(a=\) \fillintext{4},%
\par\smallskip%
\noindent\textbf{\blocktitlefont Answer}.\hypertarget{nested-tasks-3-3-2}{}\quad{}\(6\)%
\par\smallskip%
\noindent\textbf{\blocktitlefont Solution}.\hypertarget{nested-tasks-3-3-3}{}\quad{}Take the coefficient of \(x^2\) for the value of \(a\text{.}\) \emph{In this case}, \(a = {6}\text{.}\)%
\item{}\(b=\) \fillintext{4},%
\par\smallskip%
\noindent\textbf{\blocktitlefont Answer}.\hypertarget{nested-tasks-3-4-2}{}\quad{}\(-31\)%
\par\smallskip%
\noindent\textbf{\blocktitlefont Solution}.\hypertarget{nested-tasks-3-4-3}{}\quad{}Take the coefficient of \(x\) for the value of \(b\text{.}\) \emph{In this case}, \(b = {-31}\text{.}\)%
\item{}\(c=\) \fillintext{4}%
\par\smallskip%
\noindent\textbf{\blocktitlefont Answer}.\hypertarget{nested-tasks-3-5-2}{}\quad{}\(-30\)%
\par\smallskip%
\noindent\textbf{\blocktitlefont Solution}.\hypertarget{nested-tasks-3-5-3}{}\quad{}Take the constant term for the value of \(c\text{.}\) \emph{In this case}, \(c = {-30}\text{.}\)%
\end{enumerate}%
\item{}\lititle{Use the Quadratic Formula.}\par%
Using the quadratic formula, solve the equation.%
\par\smallskip%
\noindent\textbf{\blocktitlefont Answer}.\hypertarget{nested-tasks-4-3}{}\quad{}\(\left\{6,\frac{-5}{6}\right\}\)%
\par\smallskip%
\noindent\textbf{\blocktitlefont Solution}.\hypertarget{nested-tasks-4-4}{}\quad{}Recall that the quadratic formula is given in Theorem~1.2.1.%
\par
You already identified \(a = {6}\text{,}\) \(b = {-31}\text{,}\) and \(c = {-30}\text{,}\) so the results are:%
\par
%
\begin{equation*}
x = {\frac{-\left(-31\right)+\sqrt{\left(-31\right)^{2}-4\cdot 6\cdot \left(-30\right)}}{2\cdot 6}} = {6}
\end{equation*}
%
\par
or%
\par
%
\begin{equation*}
x = {\frac{-\left(-31\right)-\sqrt{\left(-31\right)^{2}-4\cdot 6\cdot \left(-30\right)}}{2\cdot 6}} = {-{\frac{5}{6}}}
\end{equation*}
%
\end{enumerate}%
This conclusion is just here for testing.%
\end{inlineexercise}%
\begin{inlineexercise}{Checkpoint}{Copy a Problem with Tasks.}{copy-with-tasks}%
We are testing copying the quadratic equation problem above (\hyperref[quadratic-equation]{Checkpoint~{\xreffont\ref{quadratic-equation}}}), since it is structured with \mono{<task>}, and we also provide a new seed.%
 \par
Consider the quadratic equation \({2x^{2}-5x-25} = 0\text{.}\)%
\begin{enumerate}[font=\bfseries,label=(\alph*),ref=\alph*]%
\item{}\lititle{Identify Coefficients.}\par%
Identify the coefficients for the quadratic equation using the standard form from Theorem~1.2.1.%
\par
\(a=\) \fillintext{4}, \(b=\) \fillintext{4}, \(c=\) \fillintext{4}%
\par\smallskip%
\noindent\textbf{\blocktitlefont Answer 1}.\hypertarget{copy-with-tasks-3-3}{}\quad{}\(2\)%
\par\smallskip%
\noindent\textbf{\blocktitlefont Answer 2}.\hypertarget{copy-with-tasks-3-4}{}\quad{}\(-5\)%
\par\smallskip%
\noindent\textbf{\blocktitlefont Answer 3}.\hypertarget{copy-with-tasks-3-5}{}\quad{}\(-25\)%
\par\smallskip%
\noindent\textbf{\blocktitlefont Solution}.\hypertarget{copy-with-tasks-3-6}{}\quad{}Take the coefficient of \(x^2\) for the value of \(a\text{,}\) the coefficient of \(x\) for \(b\text{,}\) and the constant for \(c\text{.}\) \emph{In this case}, they are \(a = {2}\text{,}\) \(b = {-5}\text{,}\) \(c = {-25}\text{.}\)%
\item{}\lititle{Use the Quadratic Formula.}\par%
Using the quadratic formula, solve the equation.%
\par\smallskip%
\noindent\textbf{\blocktitlefont Answer}.\hypertarget{copy-with-tasks-4-3}{}\quad{}\(\left\{5,\frac{-5}{2}\right\}\)%
\par\smallskip%
\noindent\textbf{\blocktitlefont Solution}.\hypertarget{copy-with-tasks-4-4}{}\quad{}Recall that the quadratic formula is given in Theorem~1.2.1.%
\par
You already identified \(a = {2}\text{,}\) \(b = {-5}\text{,}\) and \(c = {-25}\text{,}\) so the results are:%
\par
%
\begin{equation*}
x = {\frac{-\left(-5\right)+\sqrt{\left(-5\right)^{2}-4\cdot 2\cdot \left(-25\right)}}{2\cdot 2}} = {5}
\end{equation*}
%
\par
or%
\par
%
\begin{equation*}
x = {\frac{-\left(-5\right)-\sqrt{\left(-5\right)^{2}-4\cdot 2\cdot \left(-25\right)}}{2\cdot 2}} = {-{\frac{5}{2}}}
\end{equation*}
%
\end{enumerate}%
This conclusion is just here for testing.%
\end{inlineexercise}%
\end{sectionptx}
%
%
\typeout{************************************************}
\typeout{Section 1.3 Open Problem Library}
\typeout{************************************************}
%
\begin{sectionptx}{Section}{Open Problem Library}{}{Open Problem Library}{}{}{sample-ww-chapter-4-4}
The \href{http://webwork.maa.org/wiki/Open_Problem_Library}{Open Problem Library}\footnote{\nolinkurl{webwork.maa.org/wiki/Open_Problem_Library}\label{sample-ww-chapter-4-4-2-2}} (OPL) is a repository of curated WeBWorK problems. At last count it had something like 30,000 problems, searchable by topic.%
\par
Most of the problems in this sample chapter have their source authored within the same document as the narrative.  However the problems in this section are from the OPL and elsewhere on the server that is hosting the WeBWorK caclulations.  For HTML output we can connect to the server and render the problem here in an interactive fashion.  For print\slash{}PDF we use the \mono{pretext} script to connect to the server to produce PreTeXt code for the problem.  We can use the \mono{exercise.text.*} switches to control which parts of the problem (statements, solutions, hints) are incorporated.%
\par
Problems do not have to be in the OPL to reside on a server; any problem that is accessbile from the WeBWorK course which is hosting the WeBWorK computation may be used this way.%
\begin{inlineexercise}{Checkpoint}{Cylinder Volume.}{cylinder-volume}%
This problem is one that Portland Community College has donated to the Open Problem Library.%
\par
A cylinder's base's radius is \({6\ {\rm m}}\), and its height is \({10\ {\rm m}}\).%
\begin{image}{0.3}{0.4}{0.3}{}%
\includegraphics[width=\linewidth]{generated/webwork/images/webwork-32-image-1.png}
\end{image}%
%
\begin{enumerate}[label={\alph*.}]
\item{}This cylinder's volume, in terms of \(\pi\), is \fillintext{13}.%
\item{}This cylinder's volume, rounded to the hundredth place, is \fillintext{13}.%
\end{enumerate}
%
\par\smallskip%
\noindent\textbf{\blocktitlefont Answer 1}.\hypertarget{cylinder-volume-3}{}\quad{}\(360\pi\ {\rm m^{3}}\)%
\par\smallskip%
\noindent\textbf{\blocktitlefont Answer 2}.\hypertarget{cylinder-volume-4}{}\quad{}\(1130.97\ {\rm m^{3}}\)%
\par\smallskip%
\noindent\textbf{\blocktitlefont Solution}.\hypertarget{cylinder-volume-5}{}\quad{}We use \(r\) to represent the base's radius, and \(h\) to represent the cylinder's height.%
\par
A cylinder's volume formula is \(V= (\text{base area}) \cdot \text{height}\). A cylinder's base is a circle, with its area formula \(A = \pi r^{2}\).%
\par
Putting together these two formulas, we have a cylinder's volume formula:%
\par
\(\displaystyle{ V= \pi r^{2} h }\)%
\par
Throughout these computations, all quantities have units attached, and we only show them in the final step.%
\par
%
\begin{enumerate}[label={\alph*.}]
\item{}Using the volume formula, we have:%
\par
\(\displaystyle{\begin{aligned}
V \amp  = \pi r^{2} h \\
\amp  = \pi \cdot 6^{2} \cdot 10 \\
\amp  = \pi \cdot 360 \\
\amp  = 360 \pi \textrm{ m}^3
\end{aligned}}\)%
\par
Don't forget the volume unit \(\textrm{m}^3\).%
\item{}To find the decimal version, we replace \(\pi\) with its decimal value, and we have:%
\par
\(\displaystyle{\begin{aligned}[t]
V\amp  = 360 \pi \\
\amp  \approx 360 \cdot 3.14\ldots \\
\amp  \approx {1130.97\ {\rm m^{3}}}
\end{aligned}}\)%
\par
Don't forget the volume unit \(\textrm{m}^3\).%
\end{enumerate}
%
\end{inlineexercise}%
This uninteresting image, authored in \LaTeX{} syntax, is here only to make sure its automatically-assigned identifier is correct, given that the \initialism{OPL} problem just above is going to generate another image for its static representation.%
\begin{image}{0}{1}{0}{}%
\resizebox{\linewidth}{!}{%
\begin{tikzpicture}
    \draw[fill=white] (0,0) ellipse (3 and 0.45);
\end{tikzpicture}%
}%
\end{image}%
\begin{warning}{Warning}{Incompatible Problems.}{sample-ww-chapter-4-4-8}%
Not all problems that come from the OPL are compatible with PreTeXt. The reasons vary. Perhaps the problem uses older WeBWorK macros that cannot be augmented to provide PreTeXt output. Perhaps the problem nests a table within a table, which PreTeXt will not support. Perhaps there are graphics in the problem that are not sized appropriately for PreTeXt output and there's nothing you can do about it. And there are many more reasons.%
\par
Unfortunately the variety of reasons for which a problem may be incompatible means that at present, there is no way to validate a problem automatically. Using an incompatible problem might not raise any errors at any stage of the PreTeXt processing, but large chunks of the problem may be missing, especially in print output. So we recommend that you review all output modes when using an OPL problem. Chances are that if it is behaving in print, it will behave elsewhere. But that is an unproven theory at this point.%
\end{warning}
\end{sectionptx}
%
%
\typeout{************************************************}
\typeout{Section 1.4 Antidifferentiation}
\typeout{************************************************}
%
\begin{sectionptx}{Section}{Antidifferentiation}{}{Antidifferentiation}{}{}{sample-ww-chapter-4-5}
%
%
\typeout{************************************************}
\typeout{Subsection 1.4.1 A Few More Features}
\typeout{************************************************}
%
\begin{subsectionptx}{Subsection}{A Few More Features}{}{A Few More Features}{}{}{sample-ww-chapter-4-5-2}
This subsection demonstrates a few more features.%
\begin{definition}{Definition}{Antiderivative of a Function.}{definition-antiderivative}%
Suppose that \(f(x)\) and \(F(x)\) are two functions such that%
\begin{equation*}
F^\prime(x) = f(x).
\end{equation*}
Then we say \(F\) is an \terminology{antiderivative} of \(f\).%
\end{definition}
The Fundamental Theorem of Calculus in one of the high points of a course in single-variable course.%
\begin{theorem}{Theorem}{The Fundamental Theorem of Calculus.}{}{theorem-FTC}%
\index{Fundamental Theorem of Calculus}%
If \(f(x)\) is continuous, and the derivative of \(F(x)\) is \(f(x)\), then%
\begin{equation*}
\definiteintegral{a}{b}{f(x)}{x}=F(b)-F(a)
\end{equation*}
%
\end{theorem}
\begin{proof}{Proof}{}{theorem-FTC-4}
Left to the reader.%
\end{proof}
We state an equivalent version of the FTC, which is less-suited for computation, but which perhaps is a more interesting theoretical statement.%
\begin{corollary}{Corollary}{}{}{corollary-FTC-derivative}%
Suppose \(f(x)\) is a continuous function.  Then%
\begin{equation}
\frac{d}{dx}\definiteintegral{a}{x}{f(t)}{t}=f(x)\label{equation-alternate-FTC}
\end{equation}
%
\end{corollary}
\begin{proof}{Proof}{}{proof-FTC-corollary}
We simply take the indicated derivative, applying \hyperref[theorem-FTC]{Theorem~{\xreffont\ref{theorem-FTC}}} at \hyperref[equation-use-FTC]{({\xreffont\ref{equation-use-FTC}})}.%
\begin{align}
\frac{d}{dx}\definiteintegral{a}{x}{f(t)}{t}&=\frac{d}{dx}\left(F(x)-F(a)\right)\label{equation-use-FTC}\\
&=\frac{d}{dx}F(x)-\frac{d}{dx}F(a)\notag\\
&=f(x)-0 = f(x)\label{proof-FTC-corollary-1-3-3}
\end{align}
%
\end{proof}
\end{subsectionptx}
%
%
\typeout{************************************************}
\typeout{Exercises 1.4.2 WeBWorK Exercises}
\typeout{************************************************}
%
\begin{exercises-subsection}{Exercises}{WeBWorK Exercises}{}{WeBWorK Exercises}{}{}{sample-ww-chapter-4-5-3}
\begin{introduction}{}%
The first problem in this list is coming from the WeBWorK Open Problem Library. One implication of this is that we might want to provide some commentary that connects the problem to the text.  The other two ask for essay answers, which would be graded by an instructor, so in the HTML output there is no opportunity to provide an answer.%
\end{introduction}%
\begin{divisionexercise}{1}{Antiderivatives.}{}{ww-antiderivatives}%
Consult \hyperref[definition-antiderivative]{Definition~{\xreffont\ref{definition-antiderivative}}} and the \hyperref[theorem-FTC]{The Fundamental Theorem of Calculus} to assist you with the following problem.%
\par
\(\displaystyle \int_0^{5} (4 e^x+5 \sin x)\, dx\) =  \fillintext{20}%
\par
Unless the OPL has changed this problem out from under us, note the ``SOLUTION'' appearing in the solution. That is hard-coded into the OPL version of the problem. This is an example of something undesirable that may happen when using OPL problems that were not originally written with PreTeXt in mind.%
\par\smallskip%
\noindent\textbf{\blocktitlefont Answer}.\hypertarget{ww-antiderivatives-3}{}\quad{}\(593.23432548299\)%
\par\smallskip%
\noindent\textbf{\blocktitlefont Solution}.\hypertarget{ww-antiderivatives-4}{}\quad{}SOLUTION%
\par
%
\begin{equation*}
\begin{array}{rcl} \displaystyle \int_0^{5} (4 e^x+5 \sin x)\, dx \amp =\amp   
\displaystyle  4 e^x-5 \cos x  \Big]_0^{5}
\\ \amp =\amp  (4 e^{5} - 5 \cos 5) - (4 e^0 - 5 \cos0 ) 
\\ \amp =\amp  4 e^{5} - 5 \cos 5  + 1 
\end{array}
\end{equation*}
%
\end{divisionexercise}%
\begin{divisionexercise}{2}{Every Continuous Function has an Antiderivative.}{}{ww-discuss-always-antiderivative}%
WeBWorK problems can allow for open-ended essay responses that are intended to be assessed later by the instructor. For anonymous access, no text field is provided. But if this problem were used within WeBWorK as part of a homework set, users could submit an answer.%
\par
Explain how we can use Corollary~1.4.3 to say that every continuous function \emph{always} has a derivative. (And we will demonstrate here that you can use a  macro from \mono{docinfo}: \(\definiteintegral{1}{2}{\frac{1}{x}}{x}=\ln(2)\text{.}\) It will work in the WeBWorK problem, regardless of whether you are using images, MathJax, or hardcopy.)%
\end{divisionexercise}%
\begin{divisionexercise}{3}{Inverse Processes.}{}{ww-discuss-inverse-processes}%
``Differentiation and integration are inverse processes.''  Cite \alert{specific results} from this section in an explanation of how they justify this (somewhat imprecise) claim.%
\end{divisionexercise}%
\par\medskip\noindent%
\textbf{Exercise Group.}\space\space%
For the given function \(f\), find \(\indefiniteintegral{f(x)}{x}\).%
\par
Note that these common instructions are phrased in such a way that they would read well if they were applied to only one exercise at a time. That will happen if these exercises are exported as \mono{.pg} files, for example to be used in online homework from a WeBWorK server.%
\begin{exercisegroupcol}{2}
\begin{divisionexerciseegcol}{4}{}{}{sample-ww-chapter-4-5-3-6-2}%
\(f(x)=\sin(x)\)%
\par\smallskip%
\noindent\textbf{\blocktitlefont Answer}.\hypertarget{sample-ww-chapter-4-5-3-6-2-2}{}\quad{}\(-\cos\!\left(x\right)+C\)%
\end{divisionexerciseegcol}%
\begin{divisionexerciseegcol}{5}{}{}{sample-ww-chapter-4-5-3-6-3}%
\(f(x)=e^x\)%
\par\smallskip%
\noindent\textbf{\blocktitlefont Answer}.\hypertarget{sample-ww-chapter-4-5-3-6-3-2}{}\quad{}\(e^{x}+C\)%
\end{divisionexerciseegcol}%
\end{exercisegroupcol}
\par\medskip\noindent
\begin{divisionexercise}{6}{Show Your Work.}{}{ww-use-the-definition-of-the-derivative}%
Sometimes you would like a student to give a ``simple'' answer that WeBWorK can automatically assess, but you would also like the student to show their work or reasoning. Perhaps there is a particular method that you want to see the student use to find the answer. So you have a regular answer blank and also an essay blank. For practical reasons, you may wish to use the same problem on your WeBWorK server, but omit the essay part. For example, if you want to use that problem but leave out the manual grading. For this, WeBWorK has the \mono{explanation\_box} tool, demonstrated here.%
\par
Use the definition of the derivative to find \(\frac{d}{dx}x^2\text{.}\)%
\par\smallskip%
\noindent\textbf{\blocktitlefont Answer}.\hypertarget{ww-use-the-definition-of-the-derivative-3}{}\quad{}\(2x\)%
\end{divisionexercise}%
\end{exercises-subsection}
\end{sectionptx}
%
%
\typeout{************************************************}
\typeout{Section 1.5 Math Content}
\typeout{************************************************}
%
\begin{sectionptx}{Section}{Math Content}{}{Math Content}{}{}{sample-ww-chapter-4-6}
This section helps with testing aspects of math content.%
\begin{inlineexercise}{Checkpoint}{Math Elements and Alignment.}{math-and-alignment}%
In this exercise we demonstrate the allowed math elements: \mono{<m>}, \mono{<me>}, and \mono{<md>}. The last of these may have attribute \mono{@alignment} with options \mono{gather}, \mono{align}, or \mono{alignat}. The first two are used by default, depending on if you have \mono{\&} or \mono{\textbackslash{}amp} in your math.%
\par
If \(a=5\) and \(b=12\text{,}\) then%
\par
%
\begin{equation*}
a^2+b^2=13^2\text{.}
\end{equation*}
%
\par
Here, we solve an equation.%
\par
%
\begin{equation*}
\begin{aligned}
2x+1 \amp= 3\\
2x   \amp= 2\\
x    \amp= 1
\end{aligned}
\end{equation*}
%
\par
Here we have a three-way inequality to solve.%
\par
%
\begin{equation*}
\begin{gathered}
1 \lt 2x+1 \lt 3\\
0 \lt 2x   \lt 2\\
0 \lt x    \lt 1
\end{gathered}
\end{equation*}
%
\par
And here, we see a system of equations.%
\par
%
\begin{equation*}
\begin{alignedat}{3}
2x + 2y \amp {}+{} \amp z \amp {}={} \amp 10\\
y \amp {}-{} \amp  4z \amp {}={} \amp 9\\
\amp\amp   3z \amp {}={} \amp -6
\end{alignedat}
\end{equation*}
%
\end{inlineexercise}%
\begin{inlineexercise}{Checkpoint}{Intertext.}{intertext}%
With an \mono{<md>} you might have \mono{<intertext>} among the rows.%
\par
Here, we solve an equation.%
\par
%
\begin{equation*}
\begin{aligned}
2x+1 \amp= 3
\end{aligned}
\end{equation*}
Now subtract \(1\) from each side.%
\begin{equation*}
\begin{aligned}
2x   \amp= 2
\end{aligned}
\end{equation*}
Now divide by \(2\) on each side.%
\begin{equation*}
\begin{aligned}
x    \amp= 1
\end{aligned}
\end{equation*}
%
\par
We should also test when the \mono{<md>} is within a list.%
\begin{enumerate}[label={\alph*.}]
\item{}Start a list.%
\item{}~%
\par
%
\begin{equation*}
\begin{aligned}
2x+1 \amp= 3
\end{aligned}
\end{equation*}
Now subtract \(1\) from each side.%
\begin{equation*}
\begin{aligned}
2x   \amp= 2
\end{aligned}
\end{equation*}
Now divide by \(2\) on each side.%
\begin{equation*}
\begin{aligned}
x    \amp= 1
\end{aligned}
\end{equation*}
%
\item{}Still in the list?%
\end{enumerate}
%
\end{inlineexercise}%
\end{sectionptx}
%
%
\typeout{************************************************}
\typeout{Section 1.6 Multiple Choice}
\typeout{************************************************}
%
\begin{sectionptx}{Section}{Multiple Choice}{}{Multiple Choice}{}{}{sample-ww-chapter-4-7}
While free-response questions are generally preferred, sometimes the nature of a question lends itself to multiple choice.%
\begin{inlineexercise}{Checkpoint}{Drop-down\slash{}Popup.}{multiple-choice}%
Note also that the solution to this problem uses an external link.%
\par
The number \(\sqrt{2}\)\quad(\begin{itemize*}[label=$\square$,leftmargin=3em,itemjoin=\hspace{1em}]
\item{}is%
\item{}is not%
\end{itemize*})\quad
rational.%
\par\smallskip%
\noindent\textbf{\blocktitlefont Answer}.\hypertarget{multiple-choice-3}{}\quad{}\(\text{is not}\)%
\par\smallskip%
\noindent\textbf{\blocktitlefont Solution}.\hypertarget{multiple-choice-4}{}\quad{}If \(\sqrt{2}\) were rational, then \(\sqrt{2}=\frac{p}{q}\text{,}\) with \(p\) and \(q\) coprime. But then \(2q^2=p^2\text{.}\) By \href{https://en.wikipedia.org/wiki/Fundamental_theorem_of_arithmetic\#Canonical_representation_of_a_positive_integer}{the Fundamental Theorem of Arithmetic}\footnotemark{}, the power of \(2\) dividing the left side is odd, while the power of \(2\) dividing the right side is even. This is a contradiction, so \(\sqrt{2}\) is not rational.%
\end{inlineexercise}%
\footnotetext[1]{\nolinkurl{en.wikipedia.org/wiki/Fundamental_theorem_of_arithmetic\#Canonical_representation_of_a_positive_integer}\label{multiple-choice-4-1-7}}%
\begin{inlineexercise}{Checkpoint}{Choose one.}{choose-one}%
Which of the following suggest that differentiation and integration are inverse processes?%
\par
\begin{itemize}[label=$\bigcirc$,leftmargin=3em,]
\item{}The Quadratic Formula%
\item{}The Fundamental Theorem of Calculus%
\item{}The Fundamental Theorem of Arithmetic%
\item{}None of these%
\end{itemize}
%
\par\smallskip%
\noindent\textbf{\blocktitlefont Answer}.\hypertarget{choose-one-3}{}\quad{}\(\text{The Fundamental ... of Calculus}\)%
\par\smallskip%
\noindent\textbf{\blocktitlefont Solution}.\hypertarget{choose-one-4}{}\quad{}The correct answer is The Fundamental ... of Calculus.%
\end{inlineexercise}%
\begin{inlineexercise}{Checkpoint}{Choose a Subset of Options.}{options-subset}%
Select all expressions that are equivalent to \(e^{x^2 + 1/x}\text{.}\) There may be more than one correct answer.%
\par
%
\begin{itemize}[label=\textbullet]
\item{}\(\displaystyle e^{x^2} + e^{1/x}\)%
\item{}\(\displaystyle e^{ (x^3+1) / x}\)%
\item{}\(\displaystyle \dfrac{ e^{x^2} }{ e^x }\)%
\item{}\(\displaystyle e^{x^2} e^{1/x}\)%
\item{}\(\displaystyle e^{x^2} e^{x^{-1}}\)%
\item{}None of the above%
\end{itemize}
%
\par\smallskip%
\noindent\textbf{\blocktitlefont Answer}.\hypertarget{options-subset-3}{}\quad{}\(\text{Choice 2, Choice 4, Choice 5}\)%
\par\smallskip%
\noindent\textbf{\blocktitlefont Solution}.\hypertarget{options-subset-4}{}\quad{}The correct answer is Choice 2, Choice 4, Choice 5.%
\end{inlineexercise}%
\begin{inlineexercise}{Checkpoint}{Choose a Subset of Options with Automated Labeling.}{options-subset-auto-label}%
Select all expressions that are equivalent to \(e^{x^2 + 1/x}\text{.}\) There may be more than one correct answer.%
\par
%
\begin{enumerate}[label={\Alph*}]
\item{}\(\displaystyle e^{x^2} + e^{1/x}\)%
\item{}\(\displaystyle e^{x^2} e^{x^{-1}}\)%
\item{}\(\displaystyle e^{x^2} e^{1/x}\)%
\item{}\(\displaystyle e^{ (x^3+1) / x}\)%
\item{}\(\displaystyle \dfrac{ e^{x^2} }{ e^x }\)%
\item{}None of the above%
\end{enumerate}
%
\par\smallskip%
\noindent\textbf{\blocktitlefont Answer}.\hypertarget{options-subset-auto-label-3}{}\quad{}\(\text{B, C, D}\)%
\par\smallskip%
\noindent\textbf{\blocktitlefont Solution}.\hypertarget{options-subset-auto-label-4}{}\quad{}The correct answer is B, C, D.%
\end{inlineexercise}%
\begin{inlineexercise}{Checkpoint}{Choose a Subset of Options with Explicit Labeling.}{options-subset-explicit-labeling}%
Select all expressions that are equivalent to \(e^{x^2 + 1/x}\text{.}\) There may be more than one correct answer.%
\par
%
\begin{descriptionlist}
\begin{dlinarrow}{}{options-subset-explicit-labeling-2-2-1-1}%
\(\displaystyle e^{x^2} e^{1/x}\)%
\end{dlinarrow}%
\begin{dlinarrow}{}{options-subset-explicit-labeling-2-2-1-2}%
\(\displaystyle e^{x^2} + e^{1/x}\)%
\end{dlinarrow}%
\begin{dlinarrow}{}{options-subset-explicit-labeling-2-2-1-3}%
\(\displaystyle \dfrac{ e^{x^2} }{ e^x }\)%
\end{dlinarrow}%
\begin{dlinarrow}{}{options-subset-explicit-labeling-2-2-1-4}%
\(\displaystyle e^{ (x^3+1) / x}\)%
\end{dlinarrow}%
\begin{dlinarrow}{}{options-subset-explicit-labeling-2-2-1-5}%
\(\displaystyle e^{x^2} e^{x^{-1}}\)%
\end{dlinarrow}%
\begin{dlinarrow}{}{options-subset-explicit-labeling-2-2-1-6}%
None of the above%
\end{dlinarrow}%
\end{descriptionlist}
%
\par\smallskip%
\noindent\textbf{\blocktitlefont Answer}.\hypertarget{options-subset-explicit-labeling-3}{}\quad{}\(\text{TACO, SUSHI, PIZZA}\)%
\par\smallskip%
\noindent\textbf{\blocktitlefont Solution}.\hypertarget{options-subset-explicit-labeling-4}{}\quad{}The correct answer is TACO, SUSHI, PIZZA.%
\end{inlineexercise}%
\end{sectionptx}
%
%
\typeout{************************************************}
\typeout{Section 1.7 Tables}
\typeout{************************************************}
%
\begin{sectionptx}{Section}{Tables}{}{Tables}{}{}{sample-ww-chapter-4-8}
Although a WeBWorK problem written in PreTeXt source can't have a ``table'' (which would be captioned and cause trouble with free-standing PG output), they can have a ``tabular''. Tabulars can be made using the existing syntax and options that PTX offers for table-making, with some exceptions. As of summer 2017, the exceptions are that individual cells cannot have \mono{bottom} border attributes and individual columns cannot have \mono{top} border attributes. Also, while PTX table borders are generally \mono{minor}, \mono{medium}, or \mono{major}, when these borders are used in a WeBWorK problem that is rendered as a hard copy problem from within WeBWorK, the only thickness option corresponds to \mono{minor}. Several more features (like rules in general) do not carry through to static output (print or HTML-preview). However we hope this is only temporary.%
\begin{inlineexercise}{Checkpoint}{Complete this Table.}{table-completion}%
Complete this multiplication table.%
\begin{center}%
{\tabularfont%
\begin{tabular}{rCcAcA}
\(\times\)&\(6\)&\(6\)\tabularnewline\hrulethick
\(4\)&\fillintext{5}&\fillintext{5}\tabularnewline\hrulethin
\(8\)&\fillintext{5}&\fillintext{5}\tabularnewline\hrulethin
\end{tabular}
}%
\end{center}%
\par\smallskip%
\noindent\textbf{\blocktitlefont Answer 1}.\hypertarget{table-completion-3}{}\quad{}\(24\)%
\par\smallskip%
\noindent\textbf{\blocktitlefont Answer 2}.\hypertarget{table-completion-4}{}\quad{}\(24\)%
\par\smallskip%
\noindent\textbf{\blocktitlefont Answer 3}.\hypertarget{table-completion-5}{}\quad{}\(48\)%
\par\smallskip%
\noindent\textbf{\blocktitlefont Answer 4}.\hypertarget{table-completion-6}{}\quad{}\(48\)%
\par\smallskip%
\noindent\textbf{\blocktitlefont Solution}.\hypertarget{table-completion-7}{}\quad{}\begin{center}%
{\tabularfont%
\begin{tabular}{rCcAcA}
\(\times\)&\(6\)&\(6\)\tabularnewline\hrulethick
\(4\)&\(24\)&\(24\)\tabularnewline\hrulethin
\(8\)&\(48\)&\(48\)\tabularnewline\hrulethin
\end{tabular}
}%
\end{center}%
\end{inlineexercise}%
\end{sectionptx}
%
%
\typeout{************************************************}
\typeout{Section 1.8 Graphics in Exercises}
\typeout{************************************************}
%
\begin{sectionptx}{Section}{Graphics in Exercises}{}{Graphics in Exercises}{}{}{sample-ww-chapter-4-9}
It is natural for exercises to have graphics. For example, an exercise might produce a graph of some kind, and ask the reader to extract some information from that graph.%
\par
If your WeBWorK server is version 2.16 or later, WeBWorK problems can process \mono{<latex-image>} code. Here is an example.%
\begin{inlineexercise}{Checkpoint}{A static \mono{<latex-image>} graph.}{static-latex-image}%
This image is a visual proof that \(\sum_{k=1}^{n}k\) equals what?%
\begin{image}{0.25}{0.5}{0.25}{}%
\includegraphics[width=\linewidth]{generated/webwork/images/webwork-62-image-1.pdf}
\end{image}%
\par\smallskip%
\noindent\textbf{\blocktitlefont Answer}.\hypertarget{static-latex-image-3}{}\quad{}\(\mathop{\rm C}\nolimits\!\left(n+1,2\right)\hbox{ or }\frac{\left(n+1\right)n}{2}\)%
\end{inlineexercise}%
\begin{inlineexercise}{Checkpoint}{A randomized \mono{<latex-image>} graph.}{randomized-latex-image}%
These images may depend on the random seed. In this problem, the height and width of the rectangle are randomized.%
\par
Find the area of the rectangle.%
\begin{image}{0.25}{0.5}{0.25}{}%
\includegraphics[width=\linewidth]{generated/webwork/images/webwork-63-image-1.pdf}
\end{image}%
\par\smallskip%
\noindent\textbf{\blocktitlefont Answer}.\hypertarget{randomized-latex-image-3}{}\quad{}\(48\ {\rm cm^{2}}\)%
\end{inlineexercise}%
\begin{inlineexercise}{Checkpoint}{A \mono{<latex-image>} graph affected by \mono{<latex-image-preamble>}.}{latex-image-with-preamble}%
This sample chapter's \mono{<docinfo>} has a \mono{<latex-image-preamble>}. This exercise has graph styling that is affected by that.%
\par
What are the roots of this polynomial?%
\begin{image}{0.25}{0.5}{0.25}{}%
\includegraphics[width=\linewidth]{generated/webwork/images/webwork-64-image-1.pdf}
\end{image}%
\par\smallskip%
\noindent\textbf{\blocktitlefont Answer}.\hypertarget{latex-image-with-preamble-3}{}\quad{}\(-3, 0, 3\)%
\end{inlineexercise}%
\begin{inlineexercise}{Checkpoint}{Special characters.}{special-characters}%
This exercise is to test that special characters behave.%
\par
The code below has a printed dollar sign, a printed percent sign, a printed at sign, and a percent sign used as a comment marker.%
\begin{image}{0.25}{0.5}{0.25}{}%
\includegraphics[width=\linewidth]{generated/webwork/images/webwork-65-image-1.pdf}
\end{image}%
\end{inlineexercise}%
An older mechanism for creating images is supported and demonstrated here.%
\begin{inlineexercise}{Checkpoint}{Solve using a graph.}{solving-with-graph}%
The graph below is a graph of \(y=f(x)\text{.}\) Use the graph to solve the equation \(f(x)=1\text{.}\)%
\begin{image}{0.25}{0.5}{0.25}{}%
\includegraphics[width=\linewidth]{generated/webwork/images/webwork-66-image-1.png}
\end{image}%
\par\smallskip%
\noindent\textbf{\blocktitlefont Answer}.\hypertarget{solving-with-graph-3}{}\quad{}\(\left\{1\right\}\)%
\par\smallskip%
\noindent\textbf{\blocktitlefont Solution}.\hypertarget{solving-with-graph-4}{}\quad{}The graph reveals that the solution set to \(f(x)=1\) is \({\left\{1\right\}}\text{.}\)%
\begin{image}{0.25}{0.5}{0.25}{}%
\includegraphics[width=\linewidth]{generated/webwork/images/webwork-66-image-2.png}
\end{image}%
\end{inlineexercise}%
%
%
\typeout{************************************************}
\typeout{Exercises 1.8 Exercises}
\typeout{************************************************}
%
\begin{exercises-subsection-numberless}{Exercises}{Exercises}{}{Exercises}{}{}{sample-ww-chapter-4-9-10}
\par\medskip\noindent%
\textbf{Exercise Group.}\space\space%
This exercisegroup has a \mono{<latex-image>} image in its introduction. In standalone versions of the exercise, this image should be repeated.%
\begin{image}{0.25}{0.5}{0.25}{}%
\resizebox{\linewidth}{!}{%
\begin{tikzpicture}
    \draw (0,0) --++ (4,0) node[below,pos=0.5] {\(L\)} --++ (0,3) node[right,pos=0.5] {\(W\)} --++ (-4,-3) node[above,sloped,pos=0.5] {\(D\)} --++ (0,3) --++ (4,0);
\end{tikzpicture}%
}%
\end{image}%
\begin{exercisegroupcol}{2}
\begin{divisionexerciseegcol}{1}{}{}{sample-ww-chapter-4-9-10-1-2}%
Find \(D\) when \(L=4\) and \(W=3\text{.}\)%
\par\smallskip%
\noindent\textbf{\blocktitlefont Answer}.\hypertarget{sample-ww-chapter-4-9-10-1-2-2}{}\quad{}\(5\)%
\par\smallskip%
\noindent\textbf{\blocktitlefont Solution}.\hypertarget{sample-ww-chapter-4-9-10-1-2-3}{}\quad{}\begin{image}{0}{1}{0}{}%
\includegraphics[width=\linewidth]{generated/webwork/images/webwork-67-image-1.pdf}
\end{image}%
\end{divisionexerciseegcol}%
\begin{divisionexerciseegcol}{2}{}{}{sample-ww-chapter-4-9-10-1-3}%
Find \(D\) when \(L=12\) and \(W=5\text{.}\)%
\par\smallskip%
\noindent\textbf{\blocktitlefont Answer}.\hypertarget{sample-ww-chapter-4-9-10-1-3-2}{}\quad{}\(13\)%
\par\smallskip%
\noindent\textbf{\blocktitlefont Solution}.\hypertarget{sample-ww-chapter-4-9-10-1-3-3}{}\quad{}\begin{image}{0}{1}{0}{}%
\includegraphics[width=\linewidth]{generated/webwork/images/webwork-68-image-1.pdf}
\end{image}%
\end{divisionexerciseegcol}%
\end{exercisegroupcol}
\par\medskip\noindent
\end{exercises-subsection-numberless}
\end{sectionptx}
\end{chapterptx}
%
%
\typeout{************************************************}
\typeout{Chapter 2 Technical Examples}
\typeout{************************************************}
%
\begin{chapterptx}{Chapter}{Technical Examples}{}{Technical Examples}{}{}{sample-ww-chapter-5}
\renewcommand*{\chaptername}{Chapter}
%
%
\typeout{************************************************}
\typeout{Section 2.1 PGML Formatting and Verbatim Calisthenics}
\typeout{************************************************}
%
\begin{sectionptx}{Section}{PGML Formatting and Verbatim Calisthenics}{}{PGML Formatting and Verbatim Calisthenics}{}{}{sample-ww-chapter-5-2}
This section is designed to test various PGML formatting rules and verbatim content returned in answer hashes. Consult the source to see how the special characters and formatting are realized.%
\begin{inlineexercise}{Checkpoint}{PGML Formatting.}{ww-PGML-format}%
Smart double quotes: ``Life is about making an impact, not making an income.''%
\par
Smart single quotes: `Whatever the mind of man can conceive and believe, it can achieve.'%
\par
Regular apostrophes: My siblings' mother's daughter isn't my daughter's siblings' mother.%
\par
Emphasis: \emph{very important}%
\par
Alert: \alert{do not} do it%
\par
Braces: \textbraceleft{}text that looks like a set\textbraceright{}%
\begin{preformatted}
    Some pre-formatted text
                with an indented line
and an out-dented line
This should not be altered *a problem*
And this [$NDASH]* should not be an en-dash
\end{preformatted}
Here is some inline \mono{code with special characters \& < > \textquotedbl{} \textquotesingle{} \# \$ \% \textasciicircum{} \_ \{ \} \textasciitilde{} \textbackslash{} * [ ]}, and here is some%
\begin{preformatted}
single-line display code with special characters & < > " ' # $ % ^ _ { } ~ \ * [ ]
\end{preformatted}
and here is some%
\begin{preformatted}
multi-line display   
code with special characters & < > " ' # $ % ^ _ { } ~ \ * [ ]   
\end{preformatted}
Some raw characters, XML\slash{}HTML:    \& \textless{} \textgreater{} `` `%
\par
Some raw characters, \TeX{}:    \# \textdollar{} \% \textasciicircum{} \& \textunderscore{} \textbraceleft{} \textbraceright{} \textasciitilde{} \textbackslash{}%
\par
Some raw characters, PGML:        \textbackslash{} \textasteriskcentered{} \# \textbraceleft{} \textbraceright{} [ ]%
\par
Some characters that need special handling for PGML conversion to HTML or \TeX{}:%
\par
A non~breaking space (invisble where a hyphen should be)%
\par
An ndash\textendash{}right here.%
\par
An mdash\textemdash{}right here.%
\par
Some constructions in ``normal'' text, which need to be manipulated, lest they get interpreted as PGML markup:%
\par
\textgreater{}\textgreater{}This should not be a centered line \textless{}\textless{}%
\par
\textgreater{}\textgreater{}Nor a right-justified line, either%
\par
We should not get \textunderscore{}a phrase in italics\textunderscore{} in the midst of this sentence.%
\par
Brackets, unpaired; ] with content between [%
\par
Brackets, paired, in PGML constructions; [\textdollar{}NDASH]\textasteriskcentered{}, [\textunderscore{}\textunderscore{}\textunderscore{}]\textbraceleft{}\textdollar{}answer\textbraceright{}%
\par
\#\#\# This should not be a level 3 heading%
\par
===   Not a horizontal rule from three equal signs%
\par
==    Not a horizontal rule from two equal signs%
\par
===== Not a horizontal rule from five equal signs%
\par
-{}-{}-   Not a horizontal rule from three hyphens%
\par
-{}-{}    Not a horizontal rule from two hyphens%
\par
-{}-{}-{}-{}- Not a horizontal rule from five hyphens%
\par
+ Not an unordered list item%
\par
- Not an unordered list item%
\end{inlineexercise}%
Here we make MathObject String answers to see how they turn out in answer elements.%
\begin{inlineexercise}{Checkpoint}{}{sample-ww-chapter-5-2-5}%
%
\begin{enumerate}[label={\alph*.}]
\item{}Special characters used by XML, character escaping: \textless{}\textgreater{}\&'";%
\par
Now as a MathObject: \textless{}\textgreater{}\&'";%
\item{}Special characters used by LaTeX, where LaTeX \mono{\textbackslash{}text} and MathJax \mono{\textbackslash{}text} disagree: \#\%\&\textless{}\textgreater{}\textbackslash{}\textasciicircum{}\textunderscore{}\textasciigrave{}\textbar{}\textasciitilde{}%
\par
Now as a MathObject: \#\%\&\textless{}\textgreater{}\textbackslash{}\textasciicircum{}\textunderscore{}\textasciigrave{}\textbar{}\textasciitilde{}%
\item{}Special characters used by LaTeX, where LaTeX \mono{\textbackslash{}text} and MathJax \mono{\textbackslash{}text} can agree: \textdollar{}\textbraceleft{}\textbraceright{}%
\par
Now as a MathObject: \textdollar{}\textbraceleft{}\textbraceright{}%
\item{}Alphanumeric characters: ABCDEFGHIJKLMNOPQRSTUVWXYZabcdefghijklmnopqrstuvwxyz0123456789%
\par
Now as a MathObject: ABCDEFGHIJKLMNOPQRSTUVWXYZabcdefghijklmnopqrstuvwxyz0123456789%
\item{}Other characters: !()\textasteriskcentered{}+,-.\slash{}:=?@[]%
\par
Now as a MathObject: !()\textasteriskcentered{}+,-.\slash{}:=?@[]%
\end{enumerate}
%
\par
In answers, because of \TeX{} and MathJax divergence, the first two should come out in verbatim. (And so should any string containing even one of those characters.) The latter three should come out in regular text.%
\par\smallskip%
\noindent\textbf{\blocktitlefont Answer 1}.\hypertarget{sample-ww-chapter-5-2-5-2}{}\quad{}\({\verb!<>&'";!}\)%
\par\smallskip%
\noindent\textbf{\blocktitlefont Answer 2}.\hypertarget{sample-ww-chapter-5-2-5-3}{}\quad{}\({\verb!#%&<>\^_`|~!}\)%
\par\smallskip%
\noindent\textbf{\blocktitlefont Answer 3}.\hypertarget{sample-ww-chapter-5-2-5-4}{}\quad{}\({\text{\$\{\}}}\)%
\par\smallskip%
\noindent\textbf{\blocktitlefont Answer 4}.\hypertarget{sample-ww-chapter-5-2-5-5}{}\quad{}\(\text{ABCDEFGHIJKLMNOPQRSTUVWXYZabcdefghijklmnopqrstuvwxyz0123456789}\)%
\par\smallskip%
\noindent\textbf{\blocktitlefont Answer 5}.\hypertarget{sample-ww-chapter-5-2-5-6}{}\quad{}\({\text{!()*+,-./:=?@[]}}\)%
\end{inlineexercise}%
Occasionally (and probably very rarely) your PG code will include a string variable where the content is PGML syntax. If this is put into the exercise statement, it will be inserted \emph{after} PGML processing is done, and therefore it won't render as you might expect. To insert this content and also have it processed as PGML, use \mono{data=\textquotedbl{}pgml\textquotedbl{}}. (This is not in the PreTeXt schema as of 3\slash{}26\slash{}2020, and is subject to change.)%
\begin{inlineexercise}{Checkpoint}{}{sample-ww-chapter-5-2-7}%
some PGML math: [\textasciigrave{}\textbackslash{}frac\textbraceleft{}1\textbraceright{}\textbraceleft{}2\textbraceright{}+\textbackslash{}frac\textbraceleft{}3\textbraceright{}\textbraceleft{}2\textbraceright{}=2\textasciigrave{}]; and some \textasteriskcentered{}bold text\textasteriskcentered{} makes:%
\par
some PGML math: \(\frac{1}{2}+\frac{3}{2}=2\); and some \alert{bold text}%
\end{inlineexercise}%
\end{sectionptx}
%
%
\typeout{************************************************}
\typeout{Section 2.2 Subject Area Templates}
\typeout{************************************************}
%
\begin{sectionptx}{Section}{Subject Area Templates}{}{Subject Area Templates}{}{}{sample-ww-chapter-5-3}
This section samples the subject area template problems found on the WeBWorK wiki at \href{http://webwork.maa.org/wiki/SubjectAreaTemplates}{\nolinkurl{http://webwork.maa.org/wiki/SubjectAreaTemplates}}.%
\begin{inlineexercise}{Checkpoint}{Answer is a number or a function.}{number-or-function}%
%
\begin{enumerate}[label={\alph*.}]
\item{}Suppose the correct answer is \({-8}\text{.}\)%
\par
%
\begin{itemize}[label=$\blacksquare$]
\item{}\fillintext{15}%
\end{itemize}
%
\item{}Suppose the correct answer is \(\displaystyle {\frac{7-8x^{7}}{x}}\text{.}\)%
\par
%
\begin{itemize}[label=$\blacksquare$]
\item{}\fillintext{15}%
\end{itemize}
%
\end{enumerate}
%
\par\smallskip%
\noindent\textbf{\blocktitlefont Answer 1}.\hypertarget{number-or-function-3}{}\quad{}\(-8\)%
\par\smallskip%
\noindent\textbf{\blocktitlefont Answer 2}.\hypertarget{number-or-function-4}{}\quad{}\(\frac{7-8x^{7}}{x}\)%
\par\smallskip%
\noindent\textbf{\blocktitlefont Solution}.\hypertarget{number-or-function-5}{}\quad{}Solution explanation goes here.%
\end{inlineexercise}%
\begin{inlineexercise}{Checkpoint}{Answer is a function with domain issues.}{function-with-domain-issues}%
%
\begin{enumerate}[label={\alph*.}]
\item{}Suppose the correct answer is \(\displaystyle {\sqrt{x-2}}\text{.}\)%
\par
%
\begin{itemize}[label=$\blacksquare$]
\item{}\fillintext{15}%
\end{itemize}
%
\item{}Suppose the correct answer is \(\displaystyle {\ln\!\left(\left|\frac{x}{x-2}\right|\right)}\text{.}\)%
\par
%
\begin{itemize}[label=$\blacksquare$]
\item{}\fillintext{15}%
\end{itemize}
%
\end{enumerate}
%
\par\smallskip%
\noindent\textbf{\blocktitlefont Answer 1}.\hypertarget{function-with-domain-issues-3}{}\quad{}\(\sqrt{x-2}\)%
\par\smallskip%
\noindent\textbf{\blocktitlefont Answer 2}.\hypertarget{function-with-domain-issues-4}{}\quad{}\(\ln\!\left(\left|\frac{x}{x-2}\right|\right)\)%
\par\smallskip%
\noindent\textbf{\blocktitlefont Solution}.\hypertarget{function-with-domain-issues-5}{}\quad{}Solution explanation goes here.%
\end{inlineexercise}%
\begin{inlineexercise}{Checkpoint}{Multiple Choice by Popup, Radio Buttons, or Checkboxes.}{multiple-choice-popup-or-radio-buttons}%
My favorite color is\quad(\begin{itemize*}[label=$\square$,leftmargin=3em,itemjoin=\hspace{1em}]
\item{}Red%
\item{}Blue%
\item{}Green%
\end{itemize*})\quad
.%
\par
My favorite color is%
\par
\begin{itemize}[label=$\bigcirc$,leftmargin=3em,]
\item{}Red%
\item{}Blue%
\item{}Green%
\item{}None of these%
\end{itemize}
%
\par
My favorite color is%
\par
%
\begin{itemize}[label=\textbullet]
\item{}Red%
\item{}Blue%
\item{}Green%
\item{}None of these%
\end{itemize}
%
\par\smallskip%
\noindent\textbf{\blocktitlefont Answer 1}.\hypertarget{multiple-choice-popup-or-radio-buttons-3}{}\quad{}\(\text{Blue}\)%
\par\smallskip%
\noindent\textbf{\blocktitlefont Answer 2}.\hypertarget{multiple-choice-popup-or-radio-buttons-4}{}\quad{}\(\text{Blue}\)%
\par\smallskip%
\noindent\textbf{\blocktitlefont Answer 3}.\hypertarget{multiple-choice-popup-or-radio-buttons-5}{}\quad{}\(\text{Blue}\)%
\par\smallskip%
\noindent\textbf{\blocktitlefont Solution}.\hypertarget{multiple-choice-popup-or-radio-buttons-6}{}\quad{}The correct answer is Blue.%
\par
The correct answer is Blue.%
\par
The correct answer is Blue.%
\end{inlineexercise}%
\begin{inlineexercise}{Checkpoint}{}{sample-ww-chapter-5-3-6}%
There is math in each option for this question. Which expression is not a polynomial?%
\par
\begin{itemize}[label=$\bigcirc$,leftmargin=3em,]
\item{}\(\displaystyle x\)%
\item{}\(\displaystyle x^2\)%
\item{}\(\displaystyle 2^x\)%
\end{itemize}
%
\par\smallskip%
\noindent\textbf{\blocktitlefont Answer}.\hypertarget{sample-ww-chapter-5-3-6-2}{}\quad{}\(\text{Choice 3}\)%
\par\smallskip%
\noindent\textbf{\blocktitlefont Solution}.\hypertarget{sample-ww-chapter-5-3-6-3}{}\quad{}The answer is Choice 3.%
\end{inlineexercise}%
\begin{inlineexercise}{Checkpoint}{Tables.}{tables}%
A table with minimal XML source.%
\begin{center}%
{\tabularfont%
\begin{tabular}{lll}
A&B&C\tabularnewline[0pt]
D&E&F\tabularnewline[0pt]
G&H&I
\end{tabular}
}%
\end{center}%
Finish this table.%
\begin{center}%
{\tabularfont%
\begin{tabular}{CrAcClA}
\multicolumn{1}{rA}{1}&two&\(\lfloor\pi\rfloor\)\tabularnewline[0pt]
\(\text{I}+\text{I}\) \({}+\text{I}+\text{I}\)&\fillintext{5}&six\tabularnewline\crulethin{1-1}\crulethick{2-3}
\(2^3-1\)&\multicolumn{2}{cA}{VIII}\tabularnewline\hrulethick
\end{tabular}
}%
\end{center}%
\par\smallskip%
\noindent\textbf{\blocktitlefont Answer}.\hypertarget{tables-3}{}\quad{}\(5\)%
\par\smallskip%
\noindent\textbf{\blocktitlefont Solution}.\hypertarget{tables-4}{}\quad{}The missing number is 5.%
\end{inlineexercise}%
\begin{inlineexercise}{Checkpoint}{Answer Arrays.}{sample-ww-chapter-5-3-8}%
These answer blanks are all expecting some sort of answer in an array format.%
\par
\({\left(0,0\right)} = {}\)\fillintext{10} (1~\texttimes{}~2 array)%
\par
\({\left\lt 1,2\right>} = {}\)\fillintext{10} (1~\texttimes{}~2 array)%
\par
\({\left[\begin{array}{c}3\\4\\\end{array}\right]} = {}\)\fillintext{10} (2~\texttimes{}~1 array)%
\par
\({\left[\begin{array}{cc}
1 \amp 2\cr
3 \amp 4
\end{array}\right]} = {}\)\fillintext{10} (2~\texttimes{}~2 array)%
\end{inlineexercise}%
\end{sectionptx}
%
%
\typeout{************************************************}
\typeout{Section 2.3 Stress Tests}
\typeout{************************************************}
%
\begin{sectionptx}{Section}{Stress Tests}{}{Stress Tests}{}{}{sample-ww-chapter-5-4}
\begin{inlineexercise}{Checkpoint}{PTX problem source with server-generated images.}{sample-ww-chapter-5-4-2}%
\begin{enumerate}[font=\bfseries,label=(\alph*),ref=\alph*]%
\item{}\begin{image}{0.375}{0.25}{0.375}{-1.5\baselineskip}%
\includegraphics[width=\linewidth]{generated/webwork/images/webwork-79-image-1.png}
\end{image}%
\begin{image}{0.375}{0.25}{0.375}{}%
\includegraphics[width=\linewidth]{generated/webwork/images/webwork-79-image-2.png}
\end{image}%
\par\smallskip%
\noindent\textbf{\blocktitlefont Solution}.\hypertarget{sample-ww-chapter-5-4-2-2-2}{}\quad{}\begin{image}{0.375}{0.25}{0.375}{}%
\includegraphics[width=\linewidth]{generated/webwork/images/webwork-79-image-3.png}
\end{image}%
\item{}\begin{image}{0.375}{0.25}{0.375}{-1.5\baselineskip}%
\includegraphics[width=\linewidth]{generated/webwork/images/webwork-79-image-4.png}
\end{image}%
\begin{image}{0.375}{0.25}{0.375}{}%
\includegraphics[width=\linewidth]{generated/webwork/images/webwork-79-image-5.png}
\end{image}%
\begin{image}{0.375}{0.25}{0.375}{}%
\includegraphics[width=\linewidth]{generated/webwork/images/webwork-79-image-6.png}
\end{image}%
\par\smallskip%
\noindent\textbf{\blocktitlefont Hint}.\hypertarget{sample-ww-chapter-5-4-2-3-2}{}\quad{}\begin{image}{0.375}{0.25}{0.375}{}%
\includegraphics[width=\linewidth]{generated/webwork/images/webwork-79-image-7.png}
\end{image}%
\par\smallskip%
\noindent\textbf{\blocktitlefont Solution}.\hypertarget{sample-ww-chapter-5-4-2-3-3}{}\quad{}\begin{image}{0.375}{0.25}{0.375}{}%
\includegraphics[width=\linewidth]{generated/webwork/images/webwork-79-image-8.png}
\end{image}%
\end{enumerate}%
\end{inlineexercise}%
\begin{inlineexercise}{Checkpoint}{Checking Proper Indentation In Lists.}{sample-ww-chapter-5-4-3}%
One long exercise, where ordered sublists test the specification of their labels.%
\par
%
\begin{enumerate}[label={\alph*.}]
\item{}Simple item%
\item{}Simple item%
\item{}Simple item%
\end{enumerate}
%
\par
%
\begin{enumerate}[label={\alph*.}]
\item{}Simple item%
\item{}Simple item%
\item{}Simple item%
\end{enumerate}
%
\par
Text before.%
\begin{enumerate}[label={\alph*.}]
\item{}Simple item%
\item{}Simple item%
\item{}Simple item%
\end{enumerate}
%
\par
And after.%
\par
%
\begin{enumerate}[label={\alph*.}]
\item{}Structured item%
\item{}Structured item%
\item{}Structured item%
\end{enumerate}
%
\par
Text before.%
\begin{enumerate}[label={\alph*.}]
\item{}Structured item%
\item{}Structured item%
\item{}Structured item%
\end{enumerate}
%
\par
And after.%
\par
%
\begin{enumerate}[label={\alph*.}]
\item{}Structured item%
\item{}~%
\begin{enumerate}[label={\Alph*.}]
\item{}Sublist Item%
\item{}Sublist Item%
\item{}Sublist Item%
\end{enumerate}
%
\item{}Structured item%
\end{enumerate}
%
\par
Text before.%
\begin{enumerate}[label={\alph*.}]
\item{}Structured item%
\item{}Text before.%
\begin{enumerate}[label={\roman*.}]
\item{}Sublist Item%
\item{}Sublist Item%
\item{}Sublist Item%
\end{enumerate}
%
\par
And after.%
\item{}Structured item%
\end{enumerate}
%
\par
And after.%
\par
%
\begin{enumerate}[label={\alph*.}]
\item{}Structured item%
\par
Second paragraph%
\item{}~%
\begin{enumerate}[label={\Roman*.}]
\item{}Sublist Item%
\par
Second paragraph%
\item{}Sublist Item%
\par
Second paragraph%
\item{}Sublist Item%
\par
Second paragraph%
\end{enumerate}
%
\item{}Structured item%
\par
Second paragraph%
\end{enumerate}
%
\par
Text before.%
\begin{enumerate}[label={\alph*.}]
\item{}Structured item%
\par
Second paragraph%
\item{}Text before.%
\begin{enumerate}[label={\arabic*.}]
\item{}Sublist Item%
\par
Second paragraph%
\item{}Sublist Item%
\par
Second paragraph%
\item{}Sublist Item%
\par
Second paragraph%
\end{enumerate}
%
\par
And after.%
\item{}Structured item%
\par
Second paragraph%
\end{enumerate}
%
\par
And after.%
\par
%
\begin{enumerate}[label={\alph*.}]
\item{}~%
\par
%
\begin{equation*}
1+1=2
\end{equation*}
%
\par
Structured item%
\par
Second%
\par
%
\begin{equation*}
1+1=2
\end{equation*}
%
\par
paragraph%
\item{}~%
\begin{enumerate}[label={\alph*.}]
\item{}Sublist%
\par
%
\begin{equation*}
1+1=2
\end{equation*}
%
\par
Item%
\par
Second paragraph%
\par
%
\begin{equation*}
1+1=2
\end{equation*}
%
\item{}~%
\par
%
\begin{equation*}
1+1=2
\end{equation*}
%
\par
Sublist Item%
\par
Second%
\par
%
\begin{equation*}
1+1=2
\end{equation*}
%
\par
paragraph%
\item{}Sublist Item%
\par
%
\begin{equation*}
1+1=2
\end{equation*}
%
\par
%
\begin{equation*}
1+1=2
\end{equation*}
%
\par
Second paragraph%
\end{enumerate}
%
\item{}Structured%
\par
%
\begin{equation*}
1+1=2
\end{equation*}
%
\par
item%
\par
Second paragraph%
\par
%
\begin{equation*}
1+1=2
\end{equation*}
%
\end{enumerate}
%
\par
Text before.%
\begin{enumerate}[label={\alph*.}]
\item{}Structured%
\par
%
\begin{equation*}
1+1=2
\end{equation*}
%
\par
item%
\par
Second paragraph%
\par
%
\begin{equation*}
1+1=2
\end{equation*}
%
\item{}~%
\par
%
\begin{equation*}
1+1=2
\end{equation*}
%
\par
Text before.%
\begin{enumerate}[label={\Alph*.}]
\item{}Sublist%
\par
%
\begin{equation*}
1+1=2
\end{equation*}
%
\par
Item%
\par
Second paragraph%
\par
%
\begin{equation*}
1+1=2
\end{equation*}
%
\item{}~%
\par
%
\begin{equation*}
1+1=2
\end{equation*}
%
\par
Sublist Item%
\par
Second%
\par
%
\begin{equation*}
1+1=2
\end{equation*}
%
\par
paragraph%
\item{}Sublist Item%
\par
%
\begin{equation*}
1+1=2
\end{equation*}
%
\par
%
\begin{equation*}
1+1=2
\end{equation*}
%
\par
Second paragraph%
\end{enumerate}
%
\par
And%
\par
%
\begin{equation*}
1+1=2
\end{equation*}
%
\par
after.%
\item{}Structured item%
\par
%
\begin{equation*}
1+1=2
\end{equation*}
%
\par
%
\begin{equation*}
1+1=2
\end{equation*}
%
\par
Second paragraph%
\end{enumerate}
%
\par
And after.%
\par
Text before.%
\begin{enumerate}[label={\alph*.}]
\item{}Structured%
\par
%
\begin{equation*}
\begin{aligned}
1+1\amp=2\\
2\amp=2
\end{aligned}
\end{equation*}
%
\par
item%
\par
Second paragraph%
\par
%
\begin{equation*}
\begin{aligned}
1+1\amp=2\\
2\amp=2
\end{aligned}
\end{equation*}
%
\item{}~%
\par
%
\begin{equation*}
\begin{aligned}
1+1\amp=2\\
2\amp=2
\end{aligned}
\end{equation*}
%
\par
Text before.%
\begin{enumerate}[label={\roman*.}]
\item{}Sublist%
\par
%
\begin{equation*}
\begin{aligned}
1+1\amp=2\\
2\amp=2
\end{aligned}
\end{equation*}
%
\par
Item%
\par
Second paragraph%
\par
%
\begin{equation*}
\begin{aligned}
1+1\amp=2\\
2\amp=2
\end{aligned}
\end{equation*}
%
\item{}~%
\par
%
\begin{equation*}
\begin{aligned}
1+1\amp=2\\
2\amp=2
\end{aligned}
\end{equation*}
%
\par
Sublist Item%
\par
Second%
\par
%
\begin{equation*}
\begin{aligned}
1+1\amp=2\\
2\amp=2
\end{aligned}
\end{equation*}
%
\par
paragraph%
\item{}Sublist Item%
\par
%
\begin{equation*}
\begin{aligned}
1+1\amp=2\\
2\amp=2
\end{aligned}
\end{equation*}
%
\par
%
\begin{equation*}
\begin{aligned}
1+1\amp=2\\
2\amp=2
\end{aligned}
\end{equation*}
%
\par
Second paragraph%
\end{enumerate}
%
\par
And%
\par
%
\begin{equation*}
\begin{aligned}
1+1\amp=2\\
2\amp=2
\end{aligned}
\end{equation*}
%
\par
after.%
\item{}Structured item%
\par
%
\begin{equation*}
\begin{aligned}
1+1\amp=2\\
2\amp=2
\end{aligned}
\end{equation*}
%
\par
%
\begin{equation*}
\begin{aligned}
1+1\amp=2\\
2\amp=2
\end{aligned}
\end{equation*}
%
\par
Second paragraph%
\end{enumerate}
%
\par
And after.%
\end{inlineexercise}%
\begin{inlineexercise}{Checkpoint}{Checking Proper Indentation In Lists with Images and Tables.}{sample-ww-chapter-5-4-4}%
%
\begin{enumerate}[label={\alph*.}]
\item{}Structured item%
\begin{image}{0.415}{0.17}{0.415}{}%

\includegraphics[width=\linewidth]{generated/webwork/images/webwork-81-image-1.png}
\end{image}%
\item{}\begin{image}{0.415}{0.17}{0.415}{-1.5\baselineskip}%

\includegraphics[width=\linewidth]{generated/webwork/images/webwork-81-image-1.png}
\end{image}%
%
\end{enumerate}
%
\par
Structured item%
\par
%
\begin{enumerate}[label={\alph*.}]
\item{}Structured item%
\begin{image}{0.415}{0.17}{0.415}{}%

\includegraphics[width=\linewidth]{generated/webwork/images/webwork-81-image-1.png}
\end{image}%
\item{}Structured item%
\item{}Structured item%
\begin{center}%
{\tabularfont%
\begin{tabular}{ll}
a&b\tabularnewline[0pt]
c&d
\end{tabular}
}%
\end{center}%
\item{}\begin{center}%
{\tabularfont%
\begin{tabular}{CrAcClA}
\multicolumn{1}{rA}{1}&two&\(\lfloor\pi\rfloor\)\tabularnewline[0pt]
\(\text{I}+\text{I}\) \({}+\text{I}+\text{I}\)&\(5\)&six\tabularnewline\crulethin{1-1}\crulethick{2-3}
\(2^3-1\)&\multicolumn{2}{cA}{VIII}\tabularnewline\hrulethick
\end{tabular}
}%
\end{center}%
%
\end{enumerate}
%
\par
Structured item%
\par
%
\begin{enumerate}[label={\alph*.}]
\item{}Structured item%
\begin{center}%
{\tabularfont%
\begin{tabular}{ll}
a&b\tabularnewline[0pt]
c&d
\end{tabular}
}%
\end{center}%
Second paragraph%
\item{}Structured item%
\end{enumerate}
%
\end{inlineexercise}%
\begin{inlineexercise}{Checkpoint}{Deep-nested lists.}{sample-ww-chapter-5-4-5}%
Ordered list.%
\par
%
\begin{enumerate}[label={\alph*.}]
\item{}Level 1, first.%
\item{}Level 1, second.%
\begin{enumerate}[label={\roman*.}]
\item{}Level 2, first.%
\item{}Level 2, second.%
\begin{enumerate}[label={\Alph*.}]
\item{}Level 3, first.%
\item{}Level 3, second.%
\begin{enumerate}[label={\arabic*.}]
\item{}Level 4, first.%
\item{}Level 4, second.%
\item{}Level 4, third.%
\end{enumerate}
%
\item{}Level 3, third.%
\end{enumerate}
%
\item{}Level 2, third.%
\end{enumerate}
%
\item{}Level 1, third.%
\end{enumerate}
%
\par
Unordered list.%
\par
%
\begin{itemize}[label=\textbullet]
\item{}Level 1, first.%
\item{}Level 1, second.%
\begin{itemize}[label=$\circ$]
\item{}Level 2, first.%
\item{}Level 2, second.%
\begin{itemize}[label=$\blacksquare$]
\item{}Level 3, first.%
\item{}Level 3, second.%
\begin{itemize}[label=\textbullet]
\item{}Level 4, first.%
\item{}Level 4, second.%
\item{}Level 4, third.%
\end{itemize}
%
\item{}Level 3, third.%
\end{itemize}
%
\item{}Level 2, third.%
\end{itemize}
%
\item{}Level 1, third.%
\end{itemize}
%
\end{inlineexercise}%
\begin{inlineexercise}{Checkpoint}{}{sample-ww-chapter-5-4-6}%
This exercise has no content in its statement. It should throw a PTX warning during the representations build.%
\end{inlineexercise}%
\begin{inlineexercise}{Checkpoint}{}{sample-ww-chapter-5-4-7}%
This exercise has a single quote in it. A single quote is the first option for a delimiter for perl's q function which is used by extract-pg.xsl in table cells. So if working, it should move on to the next delimiter option.%
\begin{center}%
{\tabularfont%
\begin{tabular}{l}
What's up, Doc?
\end{tabular}
}%
\end{center}%
\end{inlineexercise}%
\end{sectionptx}
%
%
\typeout{************************************************}
\typeout{Section 2.4 Layout Configuration Testing}
\typeout{************************************************}
%
\begin{sectionptx}{Section}{Layout Configuration Testing}{}{Layout Configuration Testing}{}{}{sample-ww-chapter-5-5}
\begin{introduction}{}%
This section is to provide a single page demonstrating all of the combinations that may factor into a problem's layout.%
\end{introduction}%
%
%
\typeout{************************************************}
\typeout{Subsection 2.4.1 Inline Exercises}
\typeout{************************************************}
%
\begin{subsectionptx}{Subsection}{Inline Exercises}{}{Inline Exercises}{}{}{sample-ww-chapter-5-5-3}
Some ``inline'' exercises, as distinguished from the ``divisional'' exercises below.%
\begin{inlineexercise}{Checkpoint}{}{sample-ww-chapter-5-5-3-3}%
\(1+1=2\)%
\end{inlineexercise}%
\begin{inlineexercise}{Checkpoint}{Has a Title.}{sample-ww-chapter-5-5-3-4}%
\(1+1=2\)%
\end{inlineexercise}%
\begin{inlineexercise}{Checkpoint}{}{sample-ww-chapter-5-5-3-5}%
Has an introduction.%
\par
\(1+1=2\)%
\end{inlineexercise}%
\begin{inlineexercise}{Checkpoint}{Has a Title.}{sample-ww-chapter-5-5-3-6}%
Has an introduction.%
\par
\(1+1=2\)%
\end{inlineexercise}%
\begin{inlineexercise}{Checkpoint}{}{sample-ww-chapter-5-5-3-7}%
\(1+1=2\)%
\par
Has a conclusion.%
\end{inlineexercise}%
\begin{inlineexercise}{Checkpoint}{Has a Title.}{sample-ww-chapter-5-5-3-8}%
\(1+1=2\)%
\par
Has a conclusion.%
\end{inlineexercise}%
\begin{inlineexercise}{Checkpoint}{}{sample-ww-chapter-5-5-3-9}%
Has an introduction.%
\par
\(1+1=2\)%
\par
Has a conclusion.%
\end{inlineexercise}%
\begin{inlineexercise}{Checkpoint}{Has a Title.}{sample-ww-chapter-5-5-3-10}%
Has an introduction.%
\par
\(1+1=2\)%
\par
Has a conclusion.%
\end{inlineexercise}%
\end{subsectionptx}
%
%
\typeout{************************************************}
\typeout{Exercises 2.4.2 Divisional Exercises}
\typeout{************************************************}
%
\begin{exercises-subsection}{Exercises}{Divisional Exercises}{}{Divisional Exercises}{}{}{sample-ww-chapter-5-5-4}
\begin{divisionexercise}{1}{}{}{sample-ww-chapter-5-5-4-2}%
\(1+1=2\)%
\end{divisionexercise}%
\begin{divisionexercise}{2}{Has a Title.}{}{sample-ww-chapter-5-5-4-3}%
\(1+1=2\)%
\end{divisionexercise}%
\begin{divisionexercise}{3}{}{}{sample-ww-chapter-5-5-4-4}%
Has an introduction.%
\par
\(1+1=2\)%
\end{divisionexercise}%
\begin{divisionexercise}{4}{Has a Title.}{}{sample-ww-chapter-5-5-4-5}%
Has an introduction.%
\par
\(1+1=2\)%
\end{divisionexercise}%
\begin{divisionexercise}{5}{}{}{sample-ww-chapter-5-5-4-6}%
\(1+1=2\)%
\par
Has a conclusion.%
\end{divisionexercise}%
\begin{divisionexercise}{6}{Has a Title.}{}{sample-ww-chapter-5-5-4-7}%
\(1+1=2\)%
\par
Has a conclusion.%
\end{divisionexercise}%
\begin{divisionexercise}{7}{}{}{sample-ww-chapter-5-5-4-8}%
Has an introduction.%
\par
\(1+1=2\)%
\par
Has a conclusion.%
\end{divisionexercise}%
\begin{divisionexercise}{8}{Has a Title.}{}{sample-ww-chapter-5-5-4-9}%
Has an introduction.%
\par
\(1+1=2\)%
\par
Has a conclusion.%
\end{divisionexercise}%
\par\medskip\noindent%
\textbf{Exercise Group.}\space\space%
These are inside an exercisegroup.%
\begin{exercisegroup}
\begin{divisionexerciseeg}{9}{}{}{sample-ww-chapter-5-5-4-10-2}%
\(1+1=2\)%
\end{divisionexerciseeg}%
\begin{divisionexerciseeg}{10}{Has a Title.}{}{sample-ww-chapter-5-5-4-10-3}%
\(1+1=2\)%
\end{divisionexerciseeg}%
\begin{divisionexerciseeg}{11}{}{}{sample-ww-chapter-5-5-4-10-4}%
Has an introduction.%
\par
\(1+1=2\)%
\end{divisionexerciseeg}%
\begin{divisionexerciseeg}{12}{Has a Title.}{}{sample-ww-chapter-5-5-4-10-5}%
Has an introduction.%
\par
\(1+1=2\)%
\end{divisionexerciseeg}%
\begin{divisionexerciseeg}{13}{}{}{sample-ww-chapter-5-5-4-10-6}%
\(1+1=2\)%
\par
Has a conclusion.%
\end{divisionexerciseeg}%
\begin{divisionexerciseeg}{14}{Has a Title.}{}{sample-ww-chapter-5-5-4-10-7}%
\(1+1=2\)%
\par
Has a conclusion.%
\end{divisionexerciseeg}%
\begin{divisionexerciseeg}{15}{}{}{sample-ww-chapter-5-5-4-10-8}%
Has an introduction.%
\par
\(1+1=2\)%
\par
Has a conclusion.%
\end{divisionexerciseeg}%
\begin{divisionexerciseeg}{16}{Has a Title.}{}{sample-ww-chapter-5-5-4-10-9}%
Has an introduction.%
\par
\(1+1=2\)%
\par
Has a conclusion.%
\end{divisionexerciseeg}%
\end{exercisegroup}
\par\medskip\noindent
\end{exercises-subsection}
\end{sectionptx}
%
%
\typeout{************************************************}
\typeout{Section 2.5 Runestone Assignment Testing}
\typeout{************************************************}
%
\begin{sectionptx}{Section}{Runestone Assignment Testing}{}{Runestone Assignment Testing}{}{}{section-runestone-testing}
This section is specifically for testing when exercises are migrated to a Runestone Assignment page.%
%
%
\typeout{************************************************}
\typeout{Exercises 2.5 Exercises}
\typeout{************************************************}
%
\begin{exercises-subsection-numberless}{Exercises}{Exercises}{}{Exercises}{}{}{section-runestone-testing-3}
\par\medskip\noindent%
\textbf{Exercise Group.}\space\space%
This introduction should appear ahead of the exercise when it shows up in the Assignment page.%
\begin{center}%
{\tabularfont%
\begin{tabular}{ll}
\((0,0)\)&\((0,1)\)\tabularnewline[0pt]
\((1,0)\)&\((1,1)\)
\end{tabular}
}%
\end{center}%
\begin{image}{0}{1}{0}{}%
\resizebox{\linewidth}{!}{%
\begin{tikzpicture}
    \draw[fill=blue!20] (0,0) -| (1,1) -| cycle;
\end{tikzpicture}%
}%
\end{image}%
It has a table and an image too to check it all comes through.%
\begin{exercisegroup}
\begin{divisionexerciseeg}{1}{}{}{an-exercisegroup-exercise}%
What is \(1+1\text{?}\)%
\par
\fillintext{5}\textgreater{}%
\par\smallskip%
\noindent\textbf{\blocktitlefont Answer}.\hypertarget{an-exercisegroup-exercise-2}{}\quad{}\(2\)%
\end{divisionexerciseeg}%
\end{exercisegroup}
\par\medskip\noindent
\end{exercises-subsection-numberless}
\end{sectionptx}
%
%
\typeout{************************************************}
\typeout{Section 2.6 Deprecations}
\typeout{************************************************}
%
\begin{sectionptx}{Section}{Deprecations}{}{Deprecations}{}{}{section-deprecations}
This section is for testing if deprecated features are still supported. No such tests are being performed now.%
\end{sectionptx}
\end{chapterptx}
%
\appendix%
%
\clearpage\phantomsection%
\addcontentsline{toc}{part}{Appendices}%
%
%
\typeout{************************************************}
\typeout{Appendix A Hints, Answers, and Solutions}
\typeout{************************************************}
%
\begin{solutions-chapter}{Appendix}{Hints, Answers, and Solutions}{}{Hints, Answers, and Solutions}{}{}{sample-ww-chapter-6-1}
\par\medskip
\noindent\textbf{\Large{}1\space\textperiodcentered\space{}Instructive Examples\\
1.1\space\textperiodcentered\space{}Arithmetic}
\begin{inlinesolution}{Checkpoint}{1.1.1}{Adding Single-Digit Integers.}{integer-addition}%
\par\smallskip%
\noindent\textbf{\blocktitlefont Answer}.\hypertarget{integer-addition-3-back}{}\quad{}\(7\)%
\par\smallskip%
\noindent\textbf{\blocktitlefont Solution}.\hypertarget{integer-addition-4-back}{}\quad{}\(6 + 1 = {7}\text{.}\)%
\end{inlinesolution}%
\begin{inlinesolution}{Checkpoint}{1.1.2}{Declaring a Problem Seed.}{integer-addition-with-seed}%
\par\smallskip%
\noindent\textbf{\blocktitlefont Answer}.\hypertarget{integer-addition-with-seed-3-back}{}\quad{}\(13\)%
\par\smallskip%
\noindent\textbf{\blocktitlefont Solution}.\hypertarget{integer-addition-with-seed-4-back}{}\quad{}\(5 + 8 = {13}\text{.}\)%
\end{inlinesolution}%
\begin{inlinesolution}{Checkpoint}{1.1.3}{Controlling Randomness.}{integer-addition-with-control-seed}%
\par\smallskip%
\noindent\textbf{\blocktitlefont Answer}.\hypertarget{integer-addition-with-control-seed-3-back}{}\quad{}\(3\)%
\par\smallskip%
\noindent\textbf{\blocktitlefont Solution}.\hypertarget{integer-addition-with-control-seed-4-back}{}\quad{}\(1 + 2 = {3}\text{.}\)%
\end{inlinesolution}%
\begin{inlinesolution}{Checkpoint}{1.1.4}{Special Answer Checking.}{special-answer-checking}%
\par\smallskip%
\noindent\textbf{\blocktitlefont Answer}.\hypertarget{special-answer-checking-3-back}{}\quad{}\(x^{8}\)%
\par\smallskip%
\noindent\textbf{\blocktitlefont Solution}.\hypertarget{special-answer-checking-4-back}{}\quad{}We \emph{add} the exponents as follows, while including a gratuitous reference to the quadratic formula:%
\par
%
\begin{equation*}
\begin{aligned}
{x^{5}x^{3}}\amp =x^{5+3}\amp \text{Theorem 1.2.1}\\
\amp =x^{8}
\end{aligned}
\end{equation*}
%
\end{inlinesolution}%
\begin{inlinesolution}{Checkpoint}{1.1.5}{Using Hints.}{using-hints}%
\par\smallskip%
\noindent\textbf{\blocktitlefont Hint}.\hypertarget{using-hints-3-back}{}\quad{}Factor the number inside the radical.%
\par\smallskip%
\noindent\textbf{\blocktitlefont Answer}.\hypertarget{using-hints-4-back}{}\quad{}\(4\sqrt{3}\)%
\par\smallskip%
\noindent\textbf{\blocktitlefont Solution}.\hypertarget{using-hints-5-back}{}\quad{}%
\begin{equation*}
{\sqrt{48}}=\sqrt{4^2\cdot 3}=4\sqrt{3}
\end{equation*}
%
\end{inlinesolution}%
\begin{inlinesolution}{Checkpoint}{1.1.6}{No Randomization.}{no-random}%
\par\smallskip%
\noindent\textbf{\blocktitlefont Answer}.\hypertarget{no-random-3-back}{}\quad{}\(3\)%
\end{inlinesolution}%
\par\medskip
\noindent\textbf{\Large{}1.2\space\textperiodcentered\space{}The Quadratic Formula}
\begin{inlinesolution}{Checkpoint}{1.2.2}{Solving Quadratic Equations.}{quadratic-equation}%
\begin{enumerate}[font=\bfseries,label=(\alph*),ref=\alph*]%
\item[(a)]\lititle{Identify Coefficients.}\par%
\par\smallskip%
\noindent\textbf{\blocktitlefont Answer 1}.\hypertarget{quadratic-equation-3-3-back}{}\quad{}\(5\)%
\par\smallskip%
\noindent\textbf{\blocktitlefont Answer 2}.\hypertarget{quadratic-equation-3-4-back}{}\quad{}\(-6\)%
\par\smallskip%
\noindent\textbf{\blocktitlefont Answer 3}.\hypertarget{quadratic-equation-3-5-back}{}\quad{}\(-8\)%
\par\smallskip%
\noindent\textbf{\blocktitlefont Solution}.\hypertarget{quadratic-equation-3-6-back}{}\quad{}Take the coefficient of \(x^2\) for the value of \(a\text{,}\) the coefficient of \(x\) for \(b\text{,}\) and the constant for \(c\text{.}\) \emph{In this case}, they are \(a = {5}\text{,}\) \(b = {-6}\text{,}\) \(c = {-8}\text{.}\)%
\item[(b)]\lititle{Use the Quadratic Formula.}\par%
\par\smallskip%
\noindent\textbf{\blocktitlefont Answer}.\hypertarget{quadratic-equation-4-3-back}{}\quad{}\(\left\{2,\frac{-4}{5}\right\}\)%
\par\smallskip%
\noindent\textbf{\blocktitlefont Solution}.\hypertarget{quadratic-equation-4-4-back}{}\quad{}Recall that the quadratic formula is given in Theorem~1.2.1.%
\par
You already identified \(a = {5}\text{,}\) \(b = {-6}\text{,}\) and \(c = {-8}\text{,}\) so the results are:%
\par
%
\begin{equation*}
x = {\frac{-\left(-6\right)+\sqrt{\left(-6\right)^{2}-4\cdot 5\cdot \left(-8\right)}}{2\cdot 5}} = {2}
\end{equation*}
%
\par
or%
\par
%
\begin{equation*}
x = {\frac{-\left(-6\right)-\sqrt{\left(-6\right)^{2}-4\cdot 5\cdot \left(-8\right)}}{2\cdot 5}} = {-{\frac{4}{5}}}
\end{equation*}
%
\end{enumerate}%
\end{inlinesolution}%
\begin{inlinesolution}{Checkpoint}{1.2.3}{Nested tasks.}{nested-tasks}%
\begin{enumerate}[font=\bfseries,label=(\alph*),ref=\alph*]%
\item[(a)]\lititle{Identify Coefficients.}\par%
\begin{enumerate}[font=\bfseries,label=(\roman*),ref=\theenumi.\roman*]%
\item[(i)]\noindent\textbf{\blocktitlefont Answer}.\hypertarget{nested-tasks-3-3-2-back}{}\quad{}\(6\)%
\par\smallskip%
\noindent\textbf{\blocktitlefont Solution}.\hypertarget{nested-tasks-3-3-3-back}{}\quad{}Take the coefficient of \(x^2\) for the value of \(a\text{.}\) \emph{In this case}, \(a = {6}\text{.}\)%
\item[(ii)]\noindent\textbf{\blocktitlefont Answer}.\hypertarget{nested-tasks-3-4-2-back}{}\quad{}\(-31\)%
\par\smallskip%
\noindent\textbf{\blocktitlefont Solution}.\hypertarget{nested-tasks-3-4-3-back}{}\quad{}Take the coefficient of \(x\) for the value of \(b\text{.}\) \emph{In this case}, \(b = {-31}\text{.}\)%
\item[(iii)]\noindent\textbf{\blocktitlefont Answer}.\hypertarget{nested-tasks-3-5-2-back}{}\quad{}\(-30\)%
\par\smallskip%
\noindent\textbf{\blocktitlefont Solution}.\hypertarget{nested-tasks-3-5-3-back}{}\quad{}Take the constant term for the value of \(c\text{.}\) \emph{In this case}, \(c = {-30}\text{.}\)%
\end{enumerate}%
\item[(b)]\lititle{Use the Quadratic Formula.}\par%
\par\smallskip%
\noindent\textbf{\blocktitlefont Answer}.\hypertarget{nested-tasks-4-3-back}{}\quad{}\(\left\{6,\frac{-5}{6}\right\}\)%
\par\smallskip%
\noindent\textbf{\blocktitlefont Solution}.\hypertarget{nested-tasks-4-4-back}{}\quad{}Recall that the quadratic formula is given in Theorem~1.2.1.%
\par
You already identified \(a = {6}\text{,}\) \(b = {-31}\text{,}\) and \(c = {-30}\text{,}\) so the results are:%
\par
%
\begin{equation*}
x = {\frac{-\left(-31\right)+\sqrt{\left(-31\right)^{2}-4\cdot 6\cdot \left(-30\right)}}{2\cdot 6}} = {6}
\end{equation*}
%
\par
or%
\par
%
\begin{equation*}
x = {\frac{-\left(-31\right)-\sqrt{\left(-31\right)^{2}-4\cdot 6\cdot \left(-30\right)}}{2\cdot 6}} = {-{\frac{5}{6}}}
\end{equation*}
%
\end{enumerate}%
\end{inlinesolution}%
\begin{inlinesolution}{Checkpoint}{1.2.4}{Copy a Problem with Tasks.}{copy-with-tasks}%
\begin{enumerate}[font=\bfseries,label=(\alph*),ref=\alph*]%
\item[(a)]\lititle{Identify Coefficients.}\par%
\par\smallskip%
\noindent\textbf{\blocktitlefont Answer 1}.\hypertarget{copy-with-tasks-3-3-back}{}\quad{}\(2\)%
\par\smallskip%
\noindent\textbf{\blocktitlefont Answer 2}.\hypertarget{copy-with-tasks-3-4-back}{}\quad{}\(-5\)%
\par\smallskip%
\noindent\textbf{\blocktitlefont Answer 3}.\hypertarget{copy-with-tasks-3-5-back}{}\quad{}\(-25\)%
\par\smallskip%
\noindent\textbf{\blocktitlefont Solution}.\hypertarget{copy-with-tasks-3-6-back}{}\quad{}Take the coefficient of \(x^2\) for the value of \(a\text{,}\) the coefficient of \(x\) for \(b\text{,}\) and the constant for \(c\text{.}\) \emph{In this case}, they are \(a = {2}\text{,}\) \(b = {-5}\text{,}\) \(c = {-25}\text{.}\)%
\item[(b)]\lititle{Use the Quadratic Formula.}\par%
\par\smallskip%
\noindent\textbf{\blocktitlefont Answer}.\hypertarget{copy-with-tasks-4-3-back}{}\quad{}\(\left\{5,\frac{-5}{2}\right\}\)%
\par\smallskip%
\noindent\textbf{\blocktitlefont Solution}.\hypertarget{copy-with-tasks-4-4-back}{}\quad{}Recall that the quadratic formula is given in Theorem~1.2.1.%
\par
You already identified \(a = {2}\text{,}\) \(b = {-5}\text{,}\) and \(c = {-25}\text{,}\) so the results are:%
\par
%
\begin{equation*}
x = {\frac{-\left(-5\right)+\sqrt{\left(-5\right)^{2}-4\cdot 2\cdot \left(-25\right)}}{2\cdot 2}} = {5}
\end{equation*}
%
\par
or%
\par
%
\begin{equation*}
x = {\frac{-\left(-5\right)-\sqrt{\left(-5\right)^{2}-4\cdot 2\cdot \left(-25\right)}}{2\cdot 2}} = {-{\frac{5}{2}}}
\end{equation*}
%
\end{enumerate}%
\end{inlinesolution}%
\par\medskip
\noindent\textbf{\Large{}1.3\space\textperiodcentered\space{}Open Problem Library}
\begin{inlinesolution}{Checkpoint}{1.3.1}{Cylinder Volume.}{cylinder-volume}%
\par\smallskip%
\noindent\textbf{\blocktitlefont Answer 1}.\hypertarget{cylinder-volume-3-back}{}\quad{}\(360\pi\ {\rm m^{3}}\)%
\par\smallskip%
\noindent\textbf{\blocktitlefont Answer 2}.\hypertarget{cylinder-volume-4-back}{}\quad{}\(1130.97\ {\rm m^{3}}\)%
\par\smallskip%
\noindent\textbf{\blocktitlefont Solution}.\hypertarget{cylinder-volume-5-back}{}\quad{}We use \(r\) to represent the base's radius, and \(h\) to represent the cylinder's height.%
\par
A cylinder's volume formula is \(V= (\text{base area}) \cdot \text{height}\). A cylinder's base is a circle, with its area formula \(A = \pi r^{2}\).%
\par
Putting together these two formulas, we have a cylinder's volume formula:%
\par
\(\displaystyle{ V= \pi r^{2} h }\)%
\par
Throughout these computations, all quantities have units attached, and we only show them in the final step.%
\par
%
\begin{enumerate}[label={\alph*.}]
\item{}Using the volume formula, we have:%
\par
\(\displaystyle{\begin{aligned}
V \amp  = \pi r^{2} h \\
\amp  = \pi \cdot 6^{2} \cdot 10 \\
\amp  = \pi \cdot 360 \\
\amp  = 360 \pi \textrm{ m}^3
\end{aligned}}\)%
\par
Don't forget the volume unit \(\textrm{m}^3\).%
\item{}To find the decimal version, we replace \(\pi\) with its decimal value, and we have:%
\par
\(\displaystyle{\begin{aligned}[t]
V\amp  = 360 \pi \\
\amp  \approx 360 \cdot 3.14\ldots \\
\amp  \approx {1130.97\ {\rm m^{3}}}
\end{aligned}}\)%
\par
Don't forget the volume unit \(\textrm{m}^3\).%
\end{enumerate}
%
\end{inlinesolution}%
\par\medskip
\noindent\textbf{\Large{}1.4\space\textperiodcentered\space{}Antidifferentiation\\
1.4.2\space\textperiodcentered\space{}WeBWorK Exercises}
\begin{divisionsolution}{Exercise}{1.4.2.1}{Antiderivatives.}{ww-antiderivatives}%
\par\smallskip%
\noindent\textbf{\blocktitlefont Answer}.\hypertarget{ww-antiderivatives-3-back}{}\quad{}\(593.23432548299\)%
\par\smallskip%
\noindent\textbf{\blocktitlefont Solution}.\hypertarget{ww-antiderivatives-4-back}{}\quad{}SOLUTION%
\par
%
\begin{equation*}
\begin{array}{rcl} \displaystyle \int_0^{5} (4 e^x+5 \sin x)\, dx \amp =\amp   
\displaystyle  4 e^x-5 \cos x  \Big]_0^{5}
\\ \amp =\amp  (4 e^{5} - 5 \cos 5) - (4 e^0 - 5 \cos0 ) 
\\ \amp =\amp  4 e^{5} - 5 \cos 5  + 1 
\end{array}
\end{equation*}
%
\end{divisionsolution}%
\begin{exercisegroupcol}{2}
\begin{divisionsolutionegcol}{Exercise}{1.4.2.4}{}{sample-ww-chapter-4-5-3-6-2}%
\noindent\textbf{\blocktitlefont Answer}.\hypertarget{sample-ww-chapter-4-5-3-6-2-2-back}{}\quad{}\(-\cos\!\left(x\right)+C\)%
\end{divisionsolutionegcol}%
\begin{divisionsolutionegcol}{Exercise}{1.4.2.5}{}{sample-ww-chapter-4-5-3-6-3}%
\noindent\textbf{\blocktitlefont Answer}.\hypertarget{sample-ww-chapter-4-5-3-6-3-2-back}{}\quad{}\(e^{x}+C\)%
\end{divisionsolutionegcol}%
\end{exercisegroupcol}
\par\medskip\noindent
\begin{divisionsolution}{Exercise}{1.4.2.6}{Show Your Work.}{ww-use-the-definition-of-the-derivative}%
\par\smallskip%
\noindent\textbf{\blocktitlefont Answer}.\hypertarget{ww-use-the-definition-of-the-derivative-3-back}{}\quad{}\(2x\)%
\end{divisionsolution}%
\par\medskip
\noindent\textbf{\Large{}1.6\space\textperiodcentered\space{}Multiple Choice}
\begin{inlinesolution}{Checkpoint}{1.6.1}{Drop-down\slash{}Popup.}{multiple-choice}%
\par\smallskip%
\noindent\textbf{\blocktitlefont Answer}.\hypertarget{multiple-choice-3-back}{}\quad{}\(\text{is not}\)%
\par\smallskip%
\noindent\textbf{\blocktitlefont Solution}.\hypertarget{multiple-choice-4-back}{}\quad{}If \(\sqrt{2}\) were rational, then \(\sqrt{2}=\frac{p}{q}\text{,}\) with \(p\) and \(q\) coprime. But then \(2q^2=p^2\text{.}\) By \href{https://en.wikipedia.org/wiki/Fundamental_theorem_of_arithmetic\#Canonical_representation_of_a_positive_integer}{the Fundamental Theorem of Arithmetic}\footnotemark{}, the power of \(2\) dividing the left side is odd, while the power of \(2\) dividing the right side is even. This is a contradiction, so \(\sqrt{2}\) is not rational.%
\end{inlinesolution}%
\begin{inlinesolution}{Checkpoint}{1.6.2}{Choose one.}{choose-one}%
\par\smallskip%
\noindent\textbf{\blocktitlefont Answer}.\hypertarget{choose-one-3-back}{}\quad{}\(\text{The Fundamental ... of Calculus}\)%
\par\smallskip%
\noindent\textbf{\blocktitlefont Solution}.\hypertarget{choose-one-4-back}{}\quad{}The correct answer is The Fundamental ... of Calculus.%
\end{inlinesolution}%
\begin{inlinesolution}{Checkpoint}{1.6.3}{Choose a Subset of Options.}{options-subset}%
\par\smallskip%
\noindent\textbf{\blocktitlefont Answer}.\hypertarget{options-subset-3-back}{}\quad{}\(\text{Choice 2, Choice 4, Choice 5}\)%
\par\smallskip%
\noindent\textbf{\blocktitlefont Solution}.\hypertarget{options-subset-4-back}{}\quad{}The correct answer is Choice 2, Choice 4, Choice 5.%
\end{inlinesolution}%
\begin{inlinesolution}{Checkpoint}{1.6.4}{Choose a Subset of Options with Automated Labeling.}{options-subset-auto-label}%
\par\smallskip%
\noindent\textbf{\blocktitlefont Answer}.\hypertarget{options-subset-auto-label-3-back}{}\quad{}\(\text{B, C, D}\)%
\par\smallskip%
\noindent\textbf{\blocktitlefont Solution}.\hypertarget{options-subset-auto-label-4-back}{}\quad{}The correct answer is B, C, D.%
\end{inlinesolution}%
\begin{inlinesolution}{Checkpoint}{1.6.5}{Choose a Subset of Options with Explicit Labeling.}{options-subset-explicit-labeling}%
\par\smallskip%
\noindent\textbf{\blocktitlefont Answer}.\hypertarget{options-subset-explicit-labeling-3-back}{}\quad{}\(\text{TACO, SUSHI, PIZZA}\)%
\par\smallskip%
\noindent\textbf{\blocktitlefont Solution}.\hypertarget{options-subset-explicit-labeling-4-back}{}\quad{}The correct answer is TACO, SUSHI, PIZZA.%
\end{inlinesolution}%
\par\medskip
\noindent\textbf{\Large{}1.7\space\textperiodcentered\space{}Tables}
\begin{inlinesolution}{Checkpoint}{1.7.1}{Complete this Table.}{table-completion}%
\par\smallskip%
\noindent\textbf{\blocktitlefont Answer 1}.\hypertarget{table-completion-3-back}{}\quad{}\(24\)%
\par\smallskip%
\noindent\textbf{\blocktitlefont Answer 2}.\hypertarget{table-completion-4-back}{}\quad{}\(24\)%
\par\smallskip%
\noindent\textbf{\blocktitlefont Answer 3}.\hypertarget{table-completion-5-back}{}\quad{}\(48\)%
\par\smallskip%
\noindent\textbf{\blocktitlefont Answer 4}.\hypertarget{table-completion-6-back}{}\quad{}\(48\)%
\par\smallskip%
\noindent\textbf{\blocktitlefont Solution}.\hypertarget{table-completion-7-back}{}\quad{}\begin{center}%
{\tabularfont%
\begin{tabular}{rCcAcA}
\(\times\)&\(6\)&\(6\)\tabularnewline\hrulethick
\(4\)&\(24\)&\(24\)\tabularnewline\hrulethin
\(8\)&\(48\)&\(48\)\tabularnewline\hrulethin
\end{tabular}
}%
\end{center}%
\end{inlinesolution}%
\par\medskip
\noindent\textbf{\Large{}1.8\space\textperiodcentered\space{}Graphics in Exercises}
\begin{inlinesolution}{Checkpoint}{1.8.1}{A static \mono{<latex-image>} graph.}{static-latex-image}%
\par\smallskip%
\noindent\textbf{\blocktitlefont Answer}.\hypertarget{static-latex-image-3-back}{}\quad{}\(\mathop{\rm C}\nolimits\!\left(n+1,2\right)\hbox{ or }\frac{\left(n+1\right)n}{2}\)%
\end{inlinesolution}%
\begin{inlinesolution}{Checkpoint}{1.8.2}{A randomized \mono{<latex-image>} graph.}{randomized-latex-image}%
\par\smallskip%
\noindent\textbf{\blocktitlefont Answer}.\hypertarget{randomized-latex-image-3-back}{}\quad{}\(48\ {\rm cm^{2}}\)%
\end{inlinesolution}%
\begin{inlinesolution}{Checkpoint}{1.8.3}{A \mono{<latex-image>} graph affected by \mono{<latex-image-preamble>}.}{latex-image-with-preamble}%
\par\smallskip%
\noindent\textbf{\blocktitlefont Answer}.\hypertarget{latex-image-with-preamble-3-back}{}\quad{}\(-3, 0, 3\)%
\end{inlinesolution}%
\begin{inlinesolution}{Checkpoint}{1.8.5}{Solve using a graph.}{solving-with-graph}%
\par\smallskip%
\noindent\textbf{\blocktitlefont Answer}.\hypertarget{solving-with-graph-3-back}{}\quad{}\(\left\{1\right\}\)%
\par\smallskip%
\noindent\textbf{\blocktitlefont Solution}.\hypertarget{solving-with-graph-4-back}{}\quad{}The graph reveals that the solution set to \(f(x)=1\) is \({\left\{1\right\}}\text{.}\)%
\begin{image}{0.25}{0.5}{0.25}{}%
\includegraphics[width=\linewidth]{generated/webwork/images/webwork-66-image-2.png}
\end{image}%
\end{inlinesolution}%
\par\medskip
\noindent\textbf{\large{}\textperiodcentered\space{}Exercises}
\begin{exercisegroupcol}{2}
\begin{divisionsolutionegcol}{Exercise}{1.8.1}{}{sample-ww-chapter-4-9-10-1-2}%
\noindent\textbf{\blocktitlefont Answer}.\hypertarget{sample-ww-chapter-4-9-10-1-2-2-back}{}\quad{}\(5\)%
\par\smallskip%
\noindent\textbf{\blocktitlefont Solution}.\hypertarget{sample-ww-chapter-4-9-10-1-2-3-back}{}\quad{}\begin{image}{0}{1}{0}{}%
\includegraphics[width=\linewidth]{generated/webwork/images/webwork-67-image-1.pdf}
\end{image}%
\end{divisionsolutionegcol}%
\begin{divisionsolutionegcol}{Exercise}{1.8.2}{}{sample-ww-chapter-4-9-10-1-3}%
\noindent\textbf{\blocktitlefont Answer}.\hypertarget{sample-ww-chapter-4-9-10-1-3-2-back}{}\quad{}\(13\)%
\par\smallskip%
\noindent\textbf{\blocktitlefont Solution}.\hypertarget{sample-ww-chapter-4-9-10-1-3-3-back}{}\quad{}\begin{image}{0}{1}{0}{}%
\includegraphics[width=\linewidth]{generated/webwork/images/webwork-68-image-1.pdf}
\end{image}%
\end{divisionsolutionegcol}%
\end{exercisegroupcol}
\par\medskip\noindent
\par\medskip
\noindent\textbf{\Large{}2\space\textperiodcentered\space{}Technical Examples\\
2.1\space\textperiodcentered\space{}PGML Formatting and Verbatim Calisthenics}
\begin{inlinesolution}{Checkpoint}{2.1.2}{}{sample-ww-chapter-5-2-5}%
\noindent\textbf{\blocktitlefont Answer 1}.\hypertarget{sample-ww-chapter-5-2-5-2-back}{}\quad{}\({\verb!<>&'";!}\)%
\par\smallskip%
\noindent\textbf{\blocktitlefont Answer 2}.\hypertarget{sample-ww-chapter-5-2-5-3-back}{}\quad{}\({\verb!#%&<>\^_`|~!}\)%
\par\smallskip%
\noindent\textbf{\blocktitlefont Answer 3}.\hypertarget{sample-ww-chapter-5-2-5-4-back}{}\quad{}\({\text{\$\{\}}}\)%
\par\smallskip%
\noindent\textbf{\blocktitlefont Answer 4}.\hypertarget{sample-ww-chapter-5-2-5-5-back}{}\quad{}\(\text{ABCDEFGHIJKLMNOPQRSTUVWXYZabcdefghijklmnopqrstuvwxyz0123456789}\)%
\par\smallskip%
\noindent\textbf{\blocktitlefont Answer 5}.\hypertarget{sample-ww-chapter-5-2-5-6-back}{}\quad{}\({\text{!()*+,-./:=?@[]}}\)%
\end{inlinesolution}%
\par\medskip
\noindent\textbf{\Large{}2.2\space\textperiodcentered\space{}Subject Area Templates}
\begin{inlinesolution}{Checkpoint}{2.2.1}{Answer is a number or a function.}{number-or-function}%
\par\smallskip%
\noindent\textbf{\blocktitlefont Answer 1}.\hypertarget{number-or-function-3-back}{}\quad{}\(-8\)%
\par\smallskip%
\noindent\textbf{\blocktitlefont Answer 2}.\hypertarget{number-or-function-4-back}{}\quad{}\(\frac{7-8x^{7}}{x}\)%
\par\smallskip%
\noindent\textbf{\blocktitlefont Solution}.\hypertarget{number-or-function-5-back}{}\quad{}Solution explanation goes here.%
\end{inlinesolution}%
\begin{inlinesolution}{Checkpoint}{2.2.2}{Answer is a function with domain issues.}{function-with-domain-issues}%
\par\smallskip%
\noindent\textbf{\blocktitlefont Answer 1}.\hypertarget{function-with-domain-issues-3-back}{}\quad{}\(\sqrt{x-2}\)%
\par\smallskip%
\noindent\textbf{\blocktitlefont Answer 2}.\hypertarget{function-with-domain-issues-4-back}{}\quad{}\(\ln\!\left(\left|\frac{x}{x-2}\right|\right)\)%
\par\smallskip%
\noindent\textbf{\blocktitlefont Solution}.\hypertarget{function-with-domain-issues-5-back}{}\quad{}Solution explanation goes here.%
\end{inlinesolution}%
\begin{inlinesolution}{Checkpoint}{2.2.3}{Multiple Choice by Popup, Radio Buttons, or Checkboxes.}{multiple-choice-popup-or-radio-buttons}%
\par\smallskip%
\noindent\textbf{\blocktitlefont Answer 1}.\hypertarget{multiple-choice-popup-or-radio-buttons-3-back}{}\quad{}\(\text{Blue}\)%
\par\smallskip%
\noindent\textbf{\blocktitlefont Answer 2}.\hypertarget{multiple-choice-popup-or-radio-buttons-4-back}{}\quad{}\(\text{Blue}\)%
\par\smallskip%
\noindent\textbf{\blocktitlefont Answer 3}.\hypertarget{multiple-choice-popup-or-radio-buttons-5-back}{}\quad{}\(\text{Blue}\)%
\par\smallskip%
\noindent\textbf{\blocktitlefont Solution}.\hypertarget{multiple-choice-popup-or-radio-buttons-6-back}{}\quad{}The correct answer is Blue.%
\par
The correct answer is Blue.%
\par
The correct answer is Blue.%
\end{inlinesolution}%
\begin{inlinesolution}{Checkpoint}{2.2.4}{}{sample-ww-chapter-5-3-6}%
\noindent\textbf{\blocktitlefont Answer}.\hypertarget{sample-ww-chapter-5-3-6-2-back}{}\quad{}\(\text{Choice 3}\)%
\par\smallskip%
\noindent\textbf{\blocktitlefont Solution}.\hypertarget{sample-ww-chapter-5-3-6-3-back}{}\quad{}The answer is Choice 3.%
\end{inlinesolution}%
\begin{inlinesolution}{Checkpoint}{2.2.5}{Tables.}{tables}%
\par\smallskip%
\noindent\textbf{\blocktitlefont Answer}.\hypertarget{tables-3-back}{}\quad{}\(5\)%
\par\smallskip%
\noindent\textbf{\blocktitlefont Solution}.\hypertarget{tables-4-back}{}\quad{}The missing number is 5.%
\end{inlinesolution}%
\par\medskip
\noindent\textbf{\Large{}2.3\space\textperiodcentered\space{}Stress Tests}
\begin{inlinesolution}{Checkpoint}{2.3.1}{PTX problem source with server-generated images.}{sample-ww-chapter-5-4-2}%
\begin{enumerate}[font=\bfseries,label=(\alph*),ref=\alph*]%
\item[(a)]\noindent\textbf{\blocktitlefont Solution}.\hypertarget{sample-ww-chapter-5-4-2-2-2-back}{}\quad{}\begin{image}{0.375}{0.25}{0.375}{}%
\includegraphics[width=\linewidth]{generated/webwork/images/webwork-79-image-3.png}
\end{image}%
\item[(b)]\noindent\textbf{\blocktitlefont Hint}.\hypertarget{sample-ww-chapter-5-4-2-3-2-back}{}\quad{}\begin{image}{0.375}{0.25}{0.375}{}%
\includegraphics[width=\linewidth]{generated/webwork/images/webwork-79-image-7.png}
\end{image}%
\par\smallskip%
\noindent\textbf{\blocktitlefont Solution}.\hypertarget{sample-ww-chapter-5-4-2-3-3-back}{}\quad{}\begin{image}{0.375}{0.25}{0.375}{}%
\includegraphics[width=\linewidth]{generated/webwork/images/webwork-79-image-8.png}
\end{image}%
\end{enumerate}%
\end{inlinesolution}%
\par\medskip
\noindent\textbf{\Large{}2.5\space\textperiodcentered\space{}Runestone Assignment Testing}
\par\medskip
\noindent\textbf{\large{}\textperiodcentered\space{}Exercises}
\begin{exercisegroup}
\begin{divisionsolutioneg}{Exercise}{2.5.1}{}{an-exercisegroup-exercise}%
\noindent\textbf{\blocktitlefont Answer}.\hypertarget{an-exercisegroup-exercise-2-back}{}\quad{}\(2\)%
\end{divisionsolutioneg}%
\end{exercisegroup}
\par\medskip\noindent
\end{solutions-chapter}
\end{document}
